\pageoriginale

\clearpage
\SGAExposeBookmark{XVIII}{Estudio cohomológico de los pinceles de Lefschetz}
\begin{center}
{\large SGA 7 --- EXPOSÉ XVIII}\par
\vspace{2em}
{\Large\underline{Estudio cohomológico de los pinceles de Lefschetz}}\par
\vspace{1.5em}
{\large por N. Katz\footnote{Según unas notas sucintas de Grothendieck.}}
\end{center}

\section*{Sumario}

\noindent 1. La cohomología de un fibrado proyectivo. \hfill 2

\noindent 2. La cohomología de las variedades explotadas. \hfill 4

\noindent 3. La explosión asociada a un pincel. \hfill 13

\noindent 4. La cohomología de la explosión asociada a un pincel. \hfill 17

\noindent 5. Estudio cohomológico del morfismo
$\rho:\widetilde X\longrightarrow\mathbb P^1$ asociado a un pincel.
\hfill 24

\noindent\quad 5.1. El homomorfismo de restricción
$H^*(\widetilde X)\longrightarrow H^*(Y)$. \hfill 24

\noindent\quad 5.2. El ``teorema de Lefschetz vache'' (LV). \hfill 27

\noindent\quad 5.3. La condición (A). \hfill 33

\noindent\quad 5.8. El homomorfismo de Griffiths. \hfill 47

\noindent\quad 5.9. Un ejemplo en el que no se verifica la hipótesis (A).
\hfill 57

\noindent 6. Aplicaciones de la teoría local de los ciclos evanescentes
(Picard--Lefschetz). \hfill 58

\pageoriginale

\section*{Introducción}

En la presente exposición damos la estructura cohomológica elemental de los
pinceles de Lefschetz. Las secciones 1--4 recogen los hechos ``bien
conocidos'' sobre la estructura cohomológica de los fibrados proyectivos
y de las explosiones. En la sección 5 efectuamos el estudio anunciado
en el título, suponiendo que $X$ satisface el ``teorema de Lefschetz
vache'' (5.2) y que el pincel $D$ en cuestión satisface cierta condición
(A), de naturaleza local sobre $D$ (5.3). Los resultados principales
son 5.6, 5.7 y 5.8.7; este último interviene de manera esencial en la
demostración del teorema de Griffiths del Exposé XX. En la última
sección, de naturaleza claramente menos elemental, nos apoyamos en el
Exposé XIII de SGA 1, debido a Mme Raynaud, sobre el
\underline{grupo fundamental moderado}, y en la
\underline{fórmula de Picard--Lefschetz} (XVI), para demostrar:
\begin{enumerate}
\item[(1)] la condición (A) se satisface con frecuencia (6.3, 6.4);
\item[(2)] salvo cuando $\dim(X)$ es impar y la característica es $2$,
el grupo de monodromía de un pincel de Lefschetz actúa irreduciblemente
sobre la ``cohomología evanescente'' y transitivamente, salvo el signo,
sobre el conjunto de los

\pageoriginale

``ciclos evanescentes locales'' asociados, mediante la fórmula de
Picard--Lefschetz, a los diversos ``puntos críticos'' del pincel
(6.6 y 6.7).
\end{enumerate}

\section*{1. La cohomología de un fibrado proyectivo}

\subsubsection*{Teorema 1.1 (Deligne).}
\underline{Sea $f:X\longrightarrow S$ un morfismo propio, sea $\ell$ un
número primo invertible sobre $S$, sea
$u\in H^2(X,\mathbb Z_\ell(1))$ (respectivamente,
$u\in H^2(X,\mathbb Q_\ell(1))$) una clase de cohomología sobre $X$, y
sea $\mathfrak F$ un $\mathbb Z_\ell$-haz (respectivamente, un
$\mathbb Q_\ell$-haz) sobre $X$. Supongamos que existe un entero $n$
tal que, para todo entero $j$ y todo entero $i>0$, la aplicación}
\begin{equation}\tag{1.1.1}
R^{n-i}f_*\bigl(\mathfrak F(j)\bigr)
\xrightarrow{\ \wedge u^i\ }
R^{n+i}f_*\bigl(\mathfrak F(j+i)\bigr)
\end{equation}
\underline{es un isomorfismo de $\mathbb Z_\ell$-haces
(respectivamente, de $\mathbb Q_\ell$-haces) sobre $S$. Entonces, para
todo morfismo $g:S\longrightarrow T$, la sucesión espectral de Leray}
\[
E_2^{p,q}=R^pg_*R^qf_*(\mathfrak F)
\Longrightarrow R^{p+q}(gf)_*(\mathfrak F)
\]
\underline{degenera.}

La demostración se encuentra en [1, Prop.~2.4]. Obsérvese que, gracias
al teorema de cambio de base para un morfismo propio (SGA 4 XII, 5.2),
la hipótesis puede verificarse fibra por fibra.

\pageoriginale

\subsubsection*{Teorema 1.2 (cf. SGA 5 VII, 2.2.1).}
\underline{Sea $\mathcal E\longrightarrow X$ un fibrado vectorial,
localmente libre de rango $1+r$. Denotemos por
$f:\mathbb P(\mathcal E)\longrightarrow X$ el fibrado proyectivo
definido por $\mathcal E$, y por
$u\in H^2(\mathbb P(\mathcal E),\mathbb Z_\ell(1))$ la primera clase de
Chern del haz invertible canónico
$\mathcal O_{\mathbb P(\mathcal E)}(1)$. Para $\ell$ primo invertible
sobre $X$, se tiene:}
\begin{enumerate}
\item[(a)] \underline{la sucesión espectral}
\begin{equation}\tag{1.2.1}
E_2^{p,q}=H^p\bigl(X,R^qf_*\mathbb Z_\ell\bigr)
\Longrightarrow
H^{p+q}\bigl(\mathbb P(\mathcal E),\mathbb Z_\ell\bigr)
\end{equation}
\underline{degenera;}
\item[(b)] \underline{la aplicación}
\begin{equation}\tag{1.2.2}
\bigoplus_{0\leq j\leq \sisi{n}{r}}
\sisi{H^{q-2-j}\bigl(X,\mathbb Z_\ell(-j)}
      {H^{q-2j}\bigl(X,\mathbb Z_\ell(-j)\bigr)}
\longrightarrow
H^q\bigl(\mathbb P(\mathcal E),\mathbb Z_\ell\bigr),
\qquad
(\tau_{q-2j})_j\longmapsto
\sum_j f^*(\tau_{q-2j})\wedge u^j
\end{equation}
\underline{es un isomorfismo.}
\end{enumerate}
\footnote{En (1.2.2), la impresión da la cota superior $n$, el grado
$q-2-j$ y omite un paréntesis final. El rango $1+r$, la aplicación
definida y la demostración que sigue exigen $0\leq j\leq r$ y $q-2j$.}

\noindent\underline{Demostración.}
Las fibras de $f$ son espacios proyectivos de dimensión $r$, y se sabe
que
\begin{equation}\tag{1.2.3}
R^if_*\mathbb Z_\ell
\xleftarrow[\ \sim\ ]{\ \wedge^{i/2}u\ }
\begin{cases}
\mathbb Z_\ell\left(-\dfrac{i}{2}\right),
  & \text{si $i$ es par y $0\leq i\leq2r$},\\[2ex]
0, & \text{en otro caso}.
\end{cases}
\end{equation}
El teorema 1.1 se aplica a $(f,\ell,u,\mathbb Z_\ell,r)$, lo que
demuestra (a). Según (1.2.3),

\pageoriginale

\begin{equation}\tag{1.2.4}
E_2^{a,b}=
\begin{cases}
H^a\left(X,\mathbb Z_\ell\left(-\dfrac b2\right)\right),
  & \text{si $b$ es par y $0\leq b\leq2r$},\\[2ex]
0, & \text{en otro caso}.
\end{cases}
\end{equation}
Por tanto, el graduado asociado de
$H^q(\mathbb P(\mathcal E),\mathbb Z_\ell)$ para la filtración
procedente de la sucesión espectral es
\[
\bigoplus_{0\leq j\leq r}
H^{q-2j}\bigl(X,\mathbb Z_\ell(-j)\bigr).
\]
Si dotamos esta suma directa de la filtración
\[
F^i=
\bigoplus_{\substack{0\leq j\leq r\\q-2j\leq i}}
H^{q-2j}\bigl(X,\mathbb Z_\ell(-j)\bigr),
\]
la aplicación (1.2.2) es compatible con las filtraciones e induce la
identidad sobre los graduados asociados. Esto demuestra (b).

\section*{2. La cohomología de las variedades explotadas
(cf. SGA 5 VII, 8)}

\subsubsection*{2.1.}
Sean $X$ un esquema liso sobre un cuerpo $k$ y
$\Delta\subset X$ un subesquema liso en todo punto, de codimensión
$1+r$. Se sabe (EGA IV, 19.4.9) que existe un esquema
\underline{liso} $\widetilde X$, provisto de un morfismo proyectivo
\begin{equation}\tag{2.1.1}
f:\widetilde X\longrightarrow X,
\end{equation}
tal que, poniendo
\begin{equation}\tag{2.1.2}
\widetilde\Delta=f^{-1}(\Delta)
=\Delta\times_X\widetilde X,
\end{equation}
$f$ induce un isomorfismo
\begin{equation}\tag{2.1.3}
f|_{\widetilde X-\widetilde\Delta}:
\widetilde X-\widetilde\Delta
\xrightarrow{\ \sim\ }X-\Delta,
\end{equation}
y existe un $\Delta$-isomorfismo entre $\widetilde\Delta$ y

\pageoriginale

$\check{\mathbb P}(\underline N_{\Delta/X})$, el fibrado proyectivo de
rectas del fibrado normal $\underline N_{\Delta/X}$ de $\Delta$ en $X$.
Se tiene, pues, un diagrama conmutativo:
\begin{equation}\tag{2.1.4}
\begin{tikzcd}[row sep=huge,column sep=small]
\widetilde\Delta
  \arrow[rd,"f|_{\widetilde\Delta}"']
  & \simeq &
\check{\mathbb P}(\underline N_{\Delta/X})
  \arrow[ld,"\text{proyección canónica}"]\\
{} & \Delta & {}
\end{tikzcd}.
\end{equation}

\subsubsection*{2.1.5.}
En el caso especial
\[
X=\operatorname{Spec}(A),\qquad
\Delta=\operatorname{Spec}(A/I),\qquad
I=(g_0,\ldots,g_r),
\]
$\widetilde X$ se obtiene por el procedimiento de Hopf
(cf. EGA IV, 19.4.11). Se define un morfismo de $X-\Delta$ en
$\mathbb P^r$ mediante
\[
x\in X-\Delta\longmapsto
\bigl(g_0(x),\ldots,g_r(x)\bigr)\in\mathbb P^r,
\]
y se considera su gráfica $\Gamma$ inmersa en
$X\times\mathbb P^r$. Se define entonces $\widetilde X$ como la
adherencia esquemática de $\Gamma$ en $X\times\mathbb P^r$, y el
morfismo $f:\widetilde X\longrightarrow X$ se obtiene como el compuesto
\[
\widetilde X\lhook\joinrel\longrightarrow
X\times\mathbb P^r\xrightarrow{\ \operatorname{pr}_1\ }X.
\]

\subsubsection*{2.1.6.}
Análogamente (cf. de nuevo EGA IV, 19.4.11), si $X$ es proyectivo y
$\Delta$ es la intersección de $X$ con $1+r$ hipersuperficies, de
ecuaciones $G_i=0$, $i=0,\ldots,r$, se considera otra vez el morfismo
\[
X-\Delta\longrightarrow\mathbb P^r,\qquad
x\longmapsto\bigl(G_0(x),\ldots,G_r(x)\bigr),
\]
y se obtiene $\widetilde X$ como la adherencia de su gráfica en
$X\times\mathbb P^r$. Concretamente, en términos de las coordenadas

\pageoriginale

homogéneas $\lambda_0,\ldots,\lambda_r$ de $\mathbb P^r$,
$\widetilde X$ es el subesquema de $X\times\mathbb P^r$ definido por
\begin{equation}\tag{2.1.6.1}
\lambda_iG_j(X)-\lambda_jG_i(X),
\qquad 0\leq i<j\leq r.
\end{equation}
Estas ecuaciones muestran que $(x,\lambda)$ pertenece a
$\widetilde\Delta$ si y solo si $x\in\Delta$; en otras palabras,
$\check{\mathbb P}(\underline N_{\Delta/X})$ es el fibrado proyectivo
trivial
\begin{equation}\tag{2.1.6.2}
\begin{tikzcd}[column sep=large,row sep=large]
\check{\mathbb P}(\underline N_{\Delta/X})
  \arrow[rr,phantom,"\simeq"]
  \arrow[dr,"\text{proyección canónica}"']
  & {} &
\Delta\times\mathbb P^r
  \arrow[dl,"\sisi{\operatorname{pr}_i}{\operatorname{pr}_1}"]\\
{} & \Delta & {}
\end{tikzcd}.
\end{equation}
\footnote{La impresión etiqueta la proyección derecha de (2.1.6.2)
como $\operatorname{pr}_i$. Su blanco es $\Delta$, y la misma
proyección se etiqueta $\operatorname{pr}_1$ en (2.1.6.3).}
Se tiene el diagrama conmutativo
\begin{equation}\tag{2.1.6.3}
\begin{tikzcd}[column sep=large,row sep=large]
\Delta\times\mathbb P^r
  \arrow[r,"\sim"]
  \arrow[dr,"\operatorname{pr}_1"]
&
\widetilde\Delta
  \arrow[r,hook]
  \arrow[d,"f|_{\widetilde\Delta}"]
&
\widetilde X
  \arrow[r,hook]
  \arrow[d,"f"]
&
X\times\mathbb P^r
  \arrow[dl,"\operatorname{pr}_1"]
\\
{} & \Delta \arrow[r,hook] & X & {}
\end{tikzcd},
\end{equation}
y la inclusión compuesta
$\Delta\times\mathbb P^r\hookrightarrow X\times\mathbb P^r$ se obtiene
de $\Delta\hookrightarrow X$ por cambio de base.

\subsubsection*{Teorema 2.2.}
\underline{Sean $X$ liso sobre un cuerpo $k$, $\Delta$ un subesquema
liso en todo punto de codimensión $1+r$, y
$f:\widetilde X\longrightarrow X$ la explosión de $X$ a lo largo de
$\Delta$. Entonces:}

\medskip
\noindent(2.2.1) \underline{La sucesión espectral de Leray}
\[
E_2^{p,q}=H^p\bigl(X,R^qf_*\mathbb Z_\ell\bigr)
\Longrightarrow
H^{p+q}\bigl(\widetilde X,\mathbb Z_\ell\bigr)
\]
\underline{degenera.}

\pageoriginale

\medskip
\noindent(2.2.2) \underline{Existe un isomorfismo aditivo}
\[
\begin{aligned}
H^q(\widetilde X,\mathbb Z_\ell)
&\xleftarrow{\ \sim\ }
H^q(X,\mathbb Z_\ell)\oplus
\frac{H^q(\widetilde\Delta,\mathbb Z_\ell)}
     {\operatorname{Im}H^q(\Delta,\mathbb Z_\ell)}
\\
&\simeq
\left(\bigoplus_{j=1}^{r}
H^{q-2j}(\Delta,\mathbb Z_\ell(-j))\right)
\oplus H^q(X,\mathbb Z_\ell).
\end{aligned}
\]

\noindent\underline{Demostración.}
Sea
\begin{equation}\tag{2.2.3}
g:\widetilde\Delta\longrightarrow\Delta
\end{equation}
la restricción de $f$ a $\widetilde\Delta$, y sean
\begin{equation}\tag{2.2.4}
\left\{
\begin{aligned}
j&:\Delta\hookrightarrow X,\\
i&:\widetilde\Delta\hookrightarrow\widetilde X
\end{aligned}
\right.
\end{equation}
las inclusiones, de modo que se tiene el diagrama conmutativo
\begin{equation}\tag{2.2.5}
\begin{tikzcd}[column sep=huge,row sep=large]
\widetilde X \arrow[d,"f"'] &
\widetilde\Delta \arrow[l,"i"'] \arrow[d,"g"]\\
X & \Delta \arrow[l,"j"']
\end{tikzcd}.
\end{equation}
Este induce un homomorfismo entre las sucesiones espectrales de Leray
para los morfismos $f$ y $g$, con el haz $\mathbb Z_\ell$ sobre
$\widetilde X$ y $\widetilde\Delta$:
\begin{equation}\tag{2.2.6}
E_s^{p,q}(f)\xrightarrow{\ (i^*,j^*)\ }E_s^{p,q}(g).
\end{equation}
Por otra parte, como $f$ induce un isomorfismo de
$\widetilde X-\widetilde\Delta$ sobre $X-\Delta$, el teorema de cambio
de base propio da

\pageoriginale

\begin{equation}\tag{2.2.7}
R^qf_*\mathbb Z_\ell\simeq
\begin{cases}
\mathbb Z_\ell,&q=0,\\[1ex]
j_*R^qg_*\mathbb Z_\ell,&q>0.
\end{cases}
\end{equation}
Por tanto,
\begin{equation}\tag{2.2.8}
E_2^{p,q}(f)
\xrightarrow[\ \sim\ ]{\ (i^*,j^*)\ }
E_2^{p,q}(g)
\qquad(q\ne0).
\end{equation}
Demostraremos por inducción sobre $s$ que
\begin{equation}\tag{2.2.9}
E_s^{p,q}(f)
\xrightarrow[\ \sim\ ]{\ (i^*,j^*)\ }
E_s^{p,q}(g)
\qquad(q\ne0,\ s\geq2)
\end{equation}
y
\begin{equation}\tag{2.2.10}
d_s^{p,q}(f)=0\qquad(s\geq2).
\end{equation}
Admitamos por el momento la fórmula
\begin{equation}\tag{2.2.11}
d_s^{p,s-1}(f)=0\qquad(s\geq2,\ \text{para todo }p).
\end{equation}

Comencemos con $s=2$. Tenemos un diagrama conmutativo
\begin{equation}\tag{2.2.12}
\begin{tikzcd}[column sep=huge,row sep=large]
E_2^{p,q}(f)
  \arrow[r,"{(i^*,j^*)}"]
  \arrow[d,"{d_2^{p,q}(f)}"]
&
E_2^{p,q}(g)
  \arrow[d,"{0=d_2^{p,q}(g)}"]
\\
E_2^{p+2,q-1}(f)
  \arrow[r,"{(i^*,j^*)}"]
&
E_2^{p+2,q-1}(g)
\end{tikzcd}.
\end{equation}
Para $q\leq0$, se tiene $E_2^{p+2,q-1}(f)=0$, luego
$d_2^{p,q}(f)=0$. Para $q=1$, la fórmula (2.2.11) da
$d_2^{p,1}(f)=0$. Para $q>1$, las flechas horizontales de (2.2.12)
son isomorfismos, mientras que $d_2^{p,q}(g)=0$ por 1.2 y 2.1.4.
Esto demuestra (2.2.10) para $s=2$, y (2.2.9) para $s=2$ es
precisamente (2.2.8).

Supongamos conocidas (2.2.9) y (2.2.10) para cierto $s\geq2$.
La hipótesis (2.2.10) implica
$E_s^{p,q}(f)\simeq E_{s+1}^{p,q}(f)$, y por 1.2 se tiene
$E_s^{p,q}(g)\simeq E_{s+1}^{p,q}(g)$. Resulta el diagrama conmutativo

\pageoriginale

\begin{equation}\tag{2.2.13}
\begin{tikzcd}[column sep=huge,row sep=large]
E_s^{p,q}(f)
  \arrow[r,"{(i^*,j^*)}"]
  \arrow[d,"\wr"']
&
E_s^{p,q}(g)
  \arrow[d,"\wr"]
\\
E_{s+1}^{p,q}(f)
  \arrow[r,"{(i^*,j^*)}"]
&
E_{s+1}^{p,q}(g)
\end{tikzcd},
\end{equation}
que demuestra (2.2.9) para $s+1$. Para deducir (2.2.10) para
$s+1$, consideremos
\begin{equation}\tag{2.2.14}
\begin{tikzcd}[column sep=huge,row sep=large]
E_{s+1}^{p,q}(f)
  \arrow[r,"{(i^*,j^*)}"]
  \arrow[d,"{d_{s+1}^{p,q}(f)}"]
&
E_{s+1}^{p,q}(g)
  \arrow[d,"{0=d_{s+1}^{p,q}(g)}"]
\\
E_{s+1}^{p+s+1,q-s}(f)
  \arrow[r,"{(i^*,j^*)}"]
&
E_{s+1}^{p+s+1,q-s}(g)
\end{tikzcd}.
\end{equation}
Para $q\leq s-1$, se tiene
$E_{s+1}^{p+s+1,q-s}(f)=0$, luego $d_{\sisi{s}{s+1}}^{p,q}(f)=0$. Para
$q=s$, se aplica (2.2.11), y para $q>s$ las flechas horizontales de
(2.2.14) son isomorfismos y $d_{s+1}^{p,q}(g)=0$ por 1.2.
\footnote{En este punto la impresión escribe $d_s^{p,q}(f)=0$, aunque
el paso inductivo y el diagrama (2.2.14) conciernen a
$d_{s+1}^{p,q}(f)$.}

Queda verificar (2.2.11), lo que significa que el homomorfismo de borde
compuesto
\begin{equation}\tag{2.2.15}
E_2^{p,0}(f)\longrightarrow E_3^{p,0}(f)
\longrightarrow E_4^{p,0}(f)\longrightarrow\cdots
\longrightarrow E_\infty^{p,0}(f)
\end{equation}
es un isomorfismo; es decir, que
\begin{equation}\tag{2.2.16}
f^*:H^p(X,\mathbb Z_\ell)
\longrightarrow H^p(\widetilde X,\mathbb Z_\ell)
\end{equation}

\pageoriginale

es inyectivo. Pero tiene un inverso a izquierda, a saber, el
\underline{homomorfismo de Gysin} (SGA 5 IV)
\begin{equation}\tag{2.2.17}
f_*:H^p(\widetilde X,\mathbb Z_\ell)
\longrightarrow H^p(X,\mathbb Z_\ell),
\end{equation}
que satisface
\begin{equation}\tag{2.2.18}
f_*f^*=\mathrm{id}.
\end{equation}
Esto resulta de la fórmula de proyección y de que $f_*(1)=1$, pues $f$
es birracional y, por tanto, de grado uno. Así queda demostrado (2.2.1).

Para (2.2.2), al pasar a cohomología en el diagrama (2.2.5) se obtiene
\begin{equation}\tag{2.2.19}
\begin{tikzcd}[row sep=huge,column sep=huge]
H^p(\widetilde X,\mathbb Z_\ell)
  \arrow[r,"i^*"]
&
H^p(\widetilde\Delta,\mathbb Z_\ell)
\\
H^p(X,\mathbb Z_\ell)
  \arrow[r,"j^*"]
  \arrow[u,"f^*"']
&
H^p(\Delta,\mathbb Z_\ell)
  \arrow[u,"g^*"']
\end{tikzcd},
\end{equation}
y, por tanto, un morfismo
\begin{equation}\tag{2.2.20}
\frac{H^p(\widetilde X,\mathbb Z_\ell)}
     {f^*H^p(X,\mathbb Z_\ell)}
\xrightarrow{\ (i^*,j^*)\ }
\frac{H^p(\widetilde\Delta,\mathbb Z_\ell)}
     {g^*H^p(\Delta,\mathbb Z_\ell)}.
\end{equation}
Denotemos por $F^i$ las filtraciones sobre
$H^\bullet(\widetilde X,\mathbb Z_\ell)$ y
$H^\bullet(\widetilde\Delta,\mathbb Z_\ell)$ procedentes de las
sucesiones espectrales de Leray. Entonces (2.2.20) se reescribe

\pageoriginale

\begin{equation}\tag{2.2.21}
\frac{H^p(\widetilde X,\mathbb Z_\ell)}
     {F^pH^p(\widetilde X,\mathbb Z_\ell)}
\xrightarrow{\ (i^*,j^*)\ }
\frac{H^p(\widetilde\Delta,\mathbb Z_\ell)}
     {\sisi{H^pF^p}{F^pH^p}(\widetilde\Delta,\mathbb Z_\ell)}.
\end{equation}
\footnote{En el denominador derecho de (2.2.21), la impresión invierte
la notación y da $H^pF^p$; la filtración de la frase precedente y del
miembro izquierdo es $F^pH^p$.}
Por (2.2.9), la aplicación inducida por (2.2.21) sobre los graduados
asociados es un isomorfismo, por lo que (2.2.20) es un isomorfismo.
Por otra parte, (2.2.18) da una descomposición en suma directa
\begin{equation}\tag{2.2.22}
\begin{tikzcd}[column sep=7em]
H^\bullet(\widetilde X,\mathbb Z_\ell)
  \arrow[r,shift right=1.2ex,
    "f_*\oplus(1-f^*f_*)"']
&
H^\bullet(X,\mathbb Z_\ell)\oplus\operatorname{Ker} f_*
  \arrow[l,shift right=1.2ex,
    "f^*+\text{inclusión}"']
\end{tikzcd},
\end{equation}
y (1.2.2) da la descomposición
\begin{equation}\tag{2.2.23}
\frac{H^\bullet(\widetilde\Delta,\mathbb Z_\ell)}
     {g^*H^\bullet(\Delta,\mathbb Z_\ell)}
\xleftarrow{\ \sim\ }
\bigoplus_{1\leq j\leq r}
H^{\bullet-2j}(\Delta,\mathbb Z_\ell(-j)).
\end{equation}
Combinando el isomorfismo (2.2.20) y las descomposiciones (2.2.22) y
(2.2.23), se obtiene (2.2.2).

\subsubsection*{Observación 2.2.24.}
Esbocemos una segunda demostración de (2.2.2), claramente más sencilla,
válida cuando $X$ es \underline{propio} sobre $k$. Más generalmente,
sin hipótesis alguna sobre $X$, se obtiene el resultado correspondiente
(2.2.26) en cohomología con soportes propios. Las sucesiones exactas de
cohomología con soportes propios, denotada $H_c^\bullet$, para los pares
$(\widetilde X,\widetilde\Delta)$ y $(X,\Delta)$ forman un diagrama
conmutativo

\pageoriginale

\begin{center}
\resizebox{\textwidth}{!}{%
\begin{tikzcd}[ampersand replacement=\&,column sep=small,row sep=large]
(2.2.25) \& {} \arrow[r] \&
H_c^i(\widetilde X-\widetilde\Delta,\mathbb Z_\ell)
  \arrow[r] \&
H_c^i(\widetilde X,\mathbb Z_\ell)
  \arrow[r] \&
H_c^i(\widetilde\Delta,\mathbb Z_\ell)
  \arrow[r] \&
H_c^{i+1}(\widetilde X-\widetilde\Delta,\mathbb Z_\ell)
\\
(2.2.26) \& {} \arrow[r] \&
H_c^i(X-\Delta,\mathbb Z_\ell)
  \arrow[r]
  \arrow[u,"\wr" left] \&
H_c^i(X,\mathbb Z_\ell)
  \arrow[r]
  \arrow[u,hook,"f^*" right] \&
H_c^i(\Delta,\mathbb Z_\ell)
  \arrow[r]
  \arrow[u,hook,"g^*" right] \&
H_c^{i+1}(X-\Delta,\mathbb Z_\ell)
  \arrow[u,"\wr" left]
\end{tikzcd}}
\end{center}
El cociente del complejo acíclico (2.2.25) por el subcomplejo acíclico
(2.2.26) da el complejo \underline{acíclico}
\begin{equation}\tag{2.2.27}
\longrightarrow0\longrightarrow
\frac{H_c^i(\widetilde X,\mathbb Z_\ell)}
     {f^*H_c^i(X,\mathbb Z_\ell)}
\longrightarrow
\frac{H_c^i(\widetilde\Delta,\mathbb Z_\ell)}
     {g^*H_c^i(\Delta,\mathbb Z_\ell)}
\longrightarrow0\longrightarrow.
\end{equation}
Así, sin considerar sucesiones espectrales, se ve que (2.2.20) es un
isomorfismo. Los argumentos anteriores se aplican entonces y dan (2.2.2).

\subsubsection*{Corolario 2.3.}
\underline{Pongamos}
\[
b_q(X)=\dim_{\mathbb Q_\ell}H^q(\sisi{\chi}{X},\mathbb Q_\ell),
\qquad
\chi(X)=\sum_q(-1)^qb_q(X).
\]
\underline{Entonces}
\begin{equation}\tag{2.3.1}
b_q(\widetilde X)
=b_q(X)+\sum_{j=1}^{r}b_{\sisi{q-j}{q-2j}}(\Delta),
\end{equation}
\begin{equation}\tag{2.3.2}
\chi(\widetilde X)=\chi(X)+r\,\chi(\Delta).
\end{equation}

\pageoriginale

\section*{3. La explosión asociada a un pincel}

\subsection*{3.1.}
Sea $X$ proyectivo, liso e irreducible sobre un cuerpo algebraicamente
cerrado $k$, de dimensión $\sisi{u}{n}$, y sea
$D=\{H_\tau\}_{\tau\in\mathbb P^1}$ un pincel de hipersuperficies cuyo
eje corta transversalmente a $X$ en $\Delta$. Consideremos la explosión
de $X$ a lo largo de $\Delta$, que describiremos explícitamente.

En términos de ecuaciones, sean $(\lambda,\mu)$ coordenadas homogéneas
de $\mathbb P^1$, y sean $F$ y $G$ dos formas lineales tales que
$F=G=0$ defina el eje del pincel. Este se escribe
\begin{equation}\tag{3.1.1}
\lambda F-\mu G.
\end{equation}
Por definición, $\Delta$ es la intersección de $X$ con la subvariedad
lineal $G=F=0$. Por 2.1.6,
\begin{equation}\tag{3.1.2}
\widetilde X
=\{(x,(\lambda,\mu))\in X\times\mathbb P^1
       \mid \lambda F-\mu G=0\};
\end{equation}
es la adherencia esquemática, en $X\times\mathbb P^1$, de la gráfica
de la aplicación
\begin{equation}\tag{3.1.3}
\begin{tikzcd}[row sep=large,column sep=huge]
X-\Delta \arrow[r] & \mathbb P^1 \\
x \arrow[r] & (G(x),F(x)).
\end{tikzcd}
\end{equation}
La gran ventaja de $\widetilde X$ es que es proyectivo, pues está
inmerso en $X\times\mathbb P^1$, y liso, pues $\Delta$ lo es, sobre
$k$, y admite un morfismo $\rho$ hacia $\mathbb P^1$, a saber
\begin{equation}\tag{3.1.4}
\rho:\widetilde X\longrightarrow X\times\mathbb P^1
\xrightarrow{\operatorname{pr}_2}\mathbb P^1,
\end{equation}
cuyas fibras son las secciones hiperplanas de $X$ por los miembros del
pincel:
\[
\rho^{-1}(t)=X_t\overset{\mathrm{def}}{=}X\cdot H_t.
\]

\pageoriginale

\subsection*{3.2. El grado de la variedad dual: método topológico}

Supongamos además que el pincel de 3.1 es tal que
\begin{equation}\tag{3.2.1}
X_t\text{ es liso para todo }t\text{ contenido en un abierto no vacío
de }\mathbb P^1,
\end{equation}
y
\begin{equation}\tag{3.2.2}
\begin{aligned}
&\text{si $X_t$ es singular, no tiene más que singularidades cuadráticas}\\
&\text{no degeneradas (XVII 1.1).}
\end{aligned}
\end{equation}
Estas condiciones se verifican en particular para un
\underline{pincel de Lefschetz} (XVII 2.21) si
$\operatorname{char}k\ne2$ o si $X$ tiene dimensión par.

Sean $\eta$ el punto genérico de $\mathbb P^1$, $\bar\eta$ un punto
geométrico sobre él, $n=\dim X$, y $\chi$ la característica de
Euler--Poincaré con coeficientes en $\mathbb Q_\ell$. La fórmula de
Euler--Poincaré XVI 2.1, aplicada al morfismo $\sisi{S}{\rho}$, da
\begin{equation}\tag{3.2.3}
\chi(\widetilde X)
=\chi(\mathbb P^1)\,\chi(X_{\bar\eta})
 +(-1)^{\sisi{\eta}{n}}
\sum_{\substack{\text{puntos singulares}\\
                  \text{de las fibras}}}1,
\end{equation}
que, teniendo en cuenta (2.2.2), se reescribe
\begin{equation}\tag{3.2.4}
(-1)^{\sisi{\eta}{n}}
\left[\begin{array}{c}
\text{número total de}\\[3pt]
\text{puntos singulares}
\end{array}\right]
=\chi(X)+\chi(\Delta)-2\chi(X_{\bar\eta}).
\end{equation}
Obsérvese que la cantidad
$\chi(X)+\chi(\Delta)-2\chi(X_{\bar\eta})$ es independiente de la
elección particular del pincel. Depende únicamente de la inmersión de
$X$ en algún $\mathbb P^r$, e incluso solo de la clase de Chern de
$\mathcal O_X(1)$; cf. XVII, 5.7.5 y 5.7.6.

\pageoriginale

\subsubsection*{3.2.5.}
Hagamos explícito lo anterior en un caso concreto. Tomemos
$X=\mathbb P^2$ y un pincel $\{H_\tau\}$ de curvas planas de grado $d$
que satisfaga 3.1, (3.2.1) y (3.2.2). Entonces
\begin{equation}\tag{3.2.6}
\begin{array}{ll}
\chi(X)=3 & \text{(por 1.2)},\\[4pt]
\chi(X_{\bar\eta})=-d(d-3) & \text{(cf. SGA 5 VII, 7.3)},\\[4pt]
\chi(\Delta)=d^2.&
\end{array}
\end{equation}
Por tanto,
\begin{equation}\tag{3.2.7}
\text{el número total de puntos singulares es }3(d-1)^2.
\end{equation}
Tomemos el caso particular $d=3$ y consideremos la familia de curvas
elípticas, en característica distinta de $3$,
\begin{equation}\tag{3.2.8}
X^3+Y^3+Z^3+3\mu XYZ.
\end{equation}
Las cuatro fibras singulares están sobre
$\mu^3=\sisi{1}{-1}$ y $\mu=\infty$, respectivamente,\footnote{La
impresión afirma que los parámetros finitos son las raíces cúbicas de
la unidad. Con el signo positivo de (3.2.8), la factorización de Hesse
da $\mu^3=-1$; las raíces cúbicas de $+1$ corresponden al polinomio
con coeficiente $-3\mu$.}
y cada fibra singular es la unión de tres rectas:
\begin{equation}\tag{3.2.9}
\begin{tikzcd}[row sep=4.5em,column sep=4.5em]
{} & {} \arrow[dl,dash,shorten <=-9mm,shorten >=-9mm]
        \arrow[dr,dash,shorten <=-9mm,shorten >=-9mm] & {} & {} \\
{} \arrow[rr,dash,shorten <=-11mm,shorten >=-11mm] & {} & {} & {,}
\end{tikzcd}
\end{equation}
Tiene exactamente tres puntos singulares, todos ellos singularidades
cuadráticas no degeneradas. En conjunto se obtienen doce puntos
singulares, lo que verifica (3.2.7) en este ejemplo.

\pageoriginale

\subsubsection*{Proposición 3.2.10.}
\underline{Sea $X\hookrightarrow\mathbb P^r$ una inmersión de Lefschetz}
(XVII 2.3), \underline{y supongamos que $n=\dim X$ es par, o que
$\operatorname{char}k\ne2$. Entonces todo pincel de Lefschetz,
considerado como una recta de $\check{\mathbb P}^{r}$, corta a
$\check X$ en}
\begin{equation}\tag{3.2.11}
(-1)^n
\bigl[\chi(X)+\chi(\Delta)-2\chi(X_{\bar\eta})\bigr]
\end{equation}
\underline{puntos, y lo corta transversalmente} (EGA IV, 17.13.7).

\noindent\underline{Demostración.}
Por hipótesis, toda singularidad de una fibra del pincel de Lefschetz
es cuadrática no degenerada; por tanto, (3.2.4) muestra que el número
de fibras singulares del pincel $D$, cada una con exactamente un punto
singular, es (3.2.11). Ahora bien, $X_\tau$ es singular si y solo si
$\tau\in\check X\cap D$.

Para ver la transversalidad de la intersección, obsérvese que
(3.2.11) es \underline{independiente} del pincel de Lefschetz elegido.
Como el pincel genérico es de Lefschetz, (3.2.11) es el grado de
$\check X$, es decir, su grado como hipersuperficie de
$\check{\mathbb P}^{r}$, tomado por convenio igual a $0$ si
$\dim\check X\leq r-2$. Si una recta $D$ corta una hipersuperficie de
grado $d>0$ en $d$ puntos distintos $x_i$
($\sisi{i\leq i\leq d}{1\leq i\leq d}$), la intersección es
necesariamente transversal. En efecto, la multiplicidad de
intersección
\[
\mu_i=\sisi{\operatorname{rg}}{\operatorname{length}}_{x_i}
      (D\cap\check X)
\]
es al menos $1$, mientras que $\sum_{i=1}^d\mu_i=d$. Por tanto,
$\mu_i=1$ para todo $i$; equivalentemente, $D\cap\check X$ es liso,
lo cual es precisamente la transversalidad.

\pageoriginale

\subsubsection*{Observación 3.2.12.}
Para un pincel de Lefschetz $D$ del que ya se supone que corta
transversalmente a $\check X$ (condición que, en todo caso, satisface
todo $D$ suficientemente general), el resultado 3.2.10 también puede
considerarse como la conjunción de XVII 5.7.2 y del hecho
$\deg\varphi=1$ contenido en XVII 3.5. Por supuesto, el hecho probado
aquí de que \underline{todo} pincel de Lefschetz $D$ corta
transversalmente a $\check X$, mediante la fórmula de Euler--Poincaré
XVI 2.1, también puede verificarse directamente y de modo elemental,
utilizando que el eje de $D$ corta transversalmente a $X$.

\section*{4. La cohomología de la explosión asociada a un pincel}

\subsection*{4.1.}
Retomamos la situación de 3.1, resumida por el diagrama conmutativo
(cf. 2.1.6.3):
\begin{equation}\tag{4.1.1}
\resizebox{\textwidth}{!}{%
\begin{tikzcd}[ampersand replacement=\&,row sep=huge,column sep=huge]
X\times\mathbb P^1
  \& \widetilde X
      \arrow[l,hook',"f\times\rho"]
      \arrow[d,"f"]
      \arrow[dl,"\rho"']
  \& \widetilde\Delta
      \arrow[l,hook',"i"]
      \arrow[d,"g"]
      \arrow[r,"{(g,\rho_i)}","\sim"']
  \& \Delta\times\mathbb P^1 \\
\mathbb P^1 \& X \& \Delta\arrow[l,hook',"j"] \& {}
\end{tikzcd}}
\end{equation}
Comencemos por precisar (2.2.2) en este caso.

\subsection*{Proposición 4.2.}
\underline{Bajo las hipótesis de 3.1:}

\smallskip
\noindent
\textup{(4.2.1)} \underline{Los homomorfismos de
$\mathbb Z_\ell$-módulos}
\[
\begin{tikzcd}[column sep=huge]
H^\bullet(\widetilde X,\mathbb Z_\ell)
  \arrow[r,shift left=.7ex,"f_*\oplus(1-f^*f_*)"]
& H^\bullet(X,\mathbb Z_\ell)\oplus
   \operatorname{Ker}\!\left(f_*:H^\bullet(\widetilde X,\mathbb Z_\ell)
       \to H^\bullet(X,\mathbb Z_\ell)\right)
  \arrow[l,shift left=.7ex,"f^*+\operatorname{incl}"']
\end{tikzcd}
\]
\underline{son isomorfismos inversos uno de otro.}

\pageoriginale

\smallskip
\noindent
\textup{(4.2.2)} \underline{Los homomorfismos de
$\mathbb Z_\ell$-módulos}
\[
\begin{tikzcd}[column sep=huge]
H^\bullet(\widetilde X,\mathbb Z_\ell)
  \arrow[r,shift left=.7ex,
    "f_*\oplus(\sisi{}{ -}g_*i^*)"]
& H^\bullet(X,\mathbb Z_\ell)\oplus
   H^{\bullet-2}(\Delta,\mathbb Z_\ell(-1))
  \arrow[l,shift left=.7ex,"f^*+i_*g^*"']
\end{tikzcd}
\]
\underline{son isomorfismos inversos uno de otro.}
\footnote{La impresión omite el signo menos ante $g_*i^*$ en
(4.2.2), (4.2.6), (4.2.9) y (4.2.10). El lema 4.3 da
$g_*i^*i_*g^*=-\mathrm{id}$, y el propio cálculo impreso del producto
utiliza $-g_*i^*$; se restituye el signo de manera coherente.}

\smallskip
\noindent
\textup{(4.2.3)} \underline{La estructura de álgebra sobre}
\[
H^\bullet(X,\mathbb Z_\ell)\oplus
H^{\bullet-2}(\Delta,\mathbb Z_\ell(-1)),
\]
\underline{transportada mediante (4.2.2) desde el producto cup sobre}
$H^\bullet(\widetilde X,\mathbb Z_\ell)$,
\underline{se expresa, para}
\[
a,b\in H^\bullet(X,\mathbb Z_\ell),
\qquad
x,y\in H^{\bullet-2}(\Delta,\mathbb Z_\ell(-1)),
\]
\underline{por las fórmulas}
\begin{equation}\tag{4.2.3.1}
(0\oplus x)\wedge(0\oplus y)
=-j_*(xy)\oplus2xy\,c_1(\mathcal O_\Delta(1)),
\end{equation}
\begin{equation}\tag{4.2.3.2}
(a\oplus0)\wedge(b\oplus0)=ab\oplus0,
\end{equation}
\begin{equation}\tag{4.2.3.3}
(a\oplus0)\wedge(0\oplus y)=0\oplus j^*(a)y,
\end{equation}
\begin{equation}\tag{4.2.3.4}
(0\oplus x)\wedge(b\oplus0)=0\oplus xj^*(b).
\end{equation}
\underline{En particular, si}
\[
a\oplus x\in H^\bullet(X,\mathbb Z_\ell)
  \oplus H^{\bullet-2}(\Delta,\mathbb Z_\ell(-1))
\]
y
\[
b\oplus y\in H^{2n-\bullet}(X,\mathbb Z_\ell)
  \oplus H^{2n-2-\bullet}(\Delta,\mathbb Z_\ell(-1)),
\]

\pageoriginale

\underline{entonces}
\begin{equation}\tag{4.2.3.5}
(a\oplus x)\wedge(b\oplus y)=ab-j_*(xy).
\end{equation}
En otras palabras, la matriz total de intersección de
$H^\bullet(\widetilde X,\mathbb Z_\ell)$ es
\[
\left(
\begin{array}{c|c}
\begin{array}{c}
\text{matriz total de}\\[-1pt]
\text{intersección sobre $X$}
\end{array} & 0\\ \hline
0 &
-\begin{array}{c}
\text{matriz total de}\\[-1pt]
\text{intersección sobre $\Delta$}
\end{array}
\end{array}
\right).
\]

\noindent\underline{Demostración.}
Ya hemos demostrado (4.2.1); véase (2.2.22). Para (4.2.2), obsérvese
primero que, como $g:\widetilde\Delta\to\Delta$ es un fibrado en rectas
proyectivas, la sucesión
\begin{equation}\tag{4.2.4}
0\longrightarrow H^\bullet(\Delta,\mathbb Z_\ell)
\xrightarrow{g^*}H^\bullet(\widetilde\Delta,\mathbb Z_\ell)
\xrightarrow{g_*}H^{\bullet-2}(\Delta,\mathbb Z_\ell(-1))
\longrightarrow0
\end{equation}
es exacta por (1.2.2). Junto con (2.2.26), esto da el diagrama
conmutativo de sucesiones exactas
\begin{equation}\tag{4.2.5}
\resizebox{\textwidth}{!}{%
\begin{tikzcd}[ampersand replacement=\&,column sep=large,row sep=huge]
0\arrow[r]
  \&H^\bullet(X,\mathbb Z_\ell)
       \arrow[r,"f^*"]\arrow[d,"j^*" right]
  \&H^\bullet(\widetilde X,\mathbb Z_\ell)
       \arrow[r,"g_*i^*"]\arrow[d,"i^*" right]
  \&H^{\bullet-2}(\Delta,\mathbb Z_\ell(-1))
       \arrow[r]\arrow[d,equal]
  \&0\\
0\arrow[r]
  \&H^\bullet(\Delta,\mathbb Z_\ell)
       \arrow[r,"g^*"]
  \&H^\bullet(\widetilde\Delta,\mathbb Z_\ell)
       \arrow[r,"g_*"]
  \&H^{\bullet-2}(\Delta,\mathbb Z_\ell(-1))
       \arrow[r]
  \&0.
\end{tikzcd}}
\end{equation}

\pageoriginale

Teniendo en cuenta la escisión (4.2.1), se obtiene el isomorfismo
\begin{equation}\tag{4.2.6}
f_*\oplus(\sisi{}{ -}g_*i^*):
H^\bullet(\widetilde X,\mathbb Z_\ell)
\xrightarrow{\sim}
H^\bullet(X,\mathbb Z_\ell)\oplus
H^{\bullet-2}(\Delta,\mathbb Z_\ell(-1)).
\end{equation}
El transpuesto de (4.2.5), tras invertir los grados
$\bullet\mapsto2n-\bullet$, da el diagrama conmutativo de sucesiones
exactas
\begin{equation}\tag{4.2.7}
\resizebox{\textwidth}{!}{%
\begin{tikzcd}[ampersand replacement=\&,column sep=large,row sep=large]
0
  \&H^\bullet(X,\mathbb Z_\ell)\arrow[l]
  \&H^\bullet(\widetilde X,\mathbb Z_\ell)\arrow[l,"f_*"']
  \&H^{\bullet-2}(\Delta,\mathbb Z_\ell(-1))
       \arrow[l,"i_*g^*"']
  \&0\arrow[l]\\
0
  \&H^{\bullet-2}(\Delta,\mathbb Z_\ell(-2))
       \arrow[l]\arrow[u,"j_*"']
  \&H^{\bullet-2}(\widetilde\Delta,\mathbb Z_\ell(-1))
       \arrow[l,"g_*"']\arrow[u,"i_*"']
  \&H^{\bullet-2}(\Delta,\mathbb Z_\ell(-1))
       \arrow[l,"g_*"']\arrow[u,equal]
  \&0\arrow[l]
\end{tikzcd}}
\end{equation}
y, por tanto, el isomorfismo
\begin{equation}\tag{4.2.8}
H^\bullet(X,\mathbb Z_\ell)\oplus
H^{\bullet-2}(\Delta,\mathbb Z_\ell(-1))
\xrightarrow[\sim]{\ f^*+i_*g^*\ }
H^\bullet(\widetilde X,\mathbb Z_\ell).
\end{equation}
Queda demostrar que el compuesto
\begin{equation}\tag{4.2.9}
\begin{aligned}
&H^\bullet(X,\mathbb Z_\ell)\oplus
 H^{\bullet-2}(\Delta,\mathbb Z_\ell(-1))
 \xrightarrow{\ f^*+i_*g^*\ }
 H^\bullet(\widetilde X,\mathbb Z_\ell)\\[-2pt]
&\hspace{55mm}\xrightarrow{\ f_*\oplus(\sisi{}{ -}g_*i^*)\ }
 H^\bullet(X,\mathbb Z_\ell)\oplus
 H^{\bullet-2}(\Delta,\mathbb Z_\ell(-1))
\end{aligned}
\end{equation}
es la identidad. Ahora bien,
\[
f_*f^*=\mathrm{id}\quad\text{sobre }H^\bullet(X,\mathbb Z_\ell)
\quad\text{por (2.2.18),}
\]
\[
g_*i^*f^*=0\quad\text{por (4.2.5)},\qquad
f_*i_*g^*=0\quad\text{por (4.2.7)}.
\]
Por consiguiente, basta demostrar
\begin{equation}\tag{4.2.10}
\sisi{}{ -}g_*i^*i_*g^*(x)=x
\quad\text{para }x\in H^\bullet(\Delta,\mathbb Z_\ell(-1)).
\end{equation}

\pageoriginale

Para verificar (4.2.10), obsérvese primero que
$i:\widetilde\Delta\hookrightarrow\widetilde X$ es la inclusión de una
subvariedad lisa de codimensión uno. Por tanto (SGA 5 VII),
\begin{equation}\tag{4.2.11}
i^*i_*(x)=c_1(\underline N_{\widetilde\Delta/\widetilde X})\,x,
\end{equation}
y
\begin{equation}\tag{4.2.12}
i_*(x)i_*(y)
=i_*\!\left(xy\,c_1(\underline N_{\widetilde\Delta/\widetilde X})\right).
\end{equation}
Escribamos
$\xi_X=\mathcal O_X(1)$,
$\xi_\Delta=\mathcal O_\Delta(1)$ y
$\xi_{\mathbb P^1}=\mathcal O_{\mathbb P^1}(1)$.
Por EGA II, 8.1.7, si $\widetilde{\mathfrak J}$ es el ideal de
$\widetilde\Delta$ en $\widetilde X$, entonces
\begin{equation}\tag{4.2.13}
\begin{aligned}
\underline{\check N}_{\widetilde\Delta/\widetilde X}
&\simeq\widetilde{\mathfrak J}/\widetilde{\mathfrak J}^{\,2}
 \simeq\widetilde{\mathfrak J}
       \otimes_{\mathcal O_{\widetilde X}}
       \mathcal O_{\widetilde\Delta}\\
&\simeq\mathcal O_{\widetilde X}(1)
       \otimes_{\mathcal O_{\widetilde X}}
       \mathcal O_{\widetilde\Delta}
 =\mathcal O_{\widetilde\Delta}(1).
\end{aligned}
\end{equation}
Aquí $\widetilde\Delta$ se considera como
$\mathbb P(\mathfrak J/\mathfrak J^2)$, donde $\mathfrak J$ es el
ideal de $\Delta$ en $X$. Como $\Delta$ es la intersección de $X$ con
dos hiperplanos dados, existe un isomorfismo canónico
\[
\mathfrak J/\mathfrak J^2\simeq\mathcal O_\Delta^{\,2}(-1),
\]
que induce
\[
\widetilde\Delta=\mathbb P(\mathfrak J/\mathfrak J^2)
\simeq\mathbb P(\mathcal O_\Delta^{\,2}(-1))
\xrightarrow[\sim]{\alpha}
\mathbb P(\mathcal O_\Delta^{\,2})
\simeq\Delta\times\mathbb P^1.
\]
El isomorfismo canónico $\alpha$ de EGA II, 4.1.4, lleva el fibrado
tautológico en
\[
(\rho i)^*\mathcal O_{\mathbb P^1}(1)
\otimes g^*\mathcal O_\Delta(-1).
\]
En consecuencia, (4.2.13) se convierte en
\begin{equation}\tag{4.2.14}
\underline{\check N}_{\widetilde\Delta/\widetilde X}
\simeq(\rho i)^*\mathcal O_{\mathbb P^1}(1)
       \otimes g^*\mathcal O_\Delta(-1).
\end{equation}

\pageoriginale

\subsection*{Lema 4.3.}
\underline{Se tiene}
\[
c_1(\underline N_{\widetilde\Delta/\widetilde X})
=g^*c_1(\xi_\Delta)-(\rho i)^*c_1(\xi_{\mathbb P^1}).
\]
Volvamos a la demostración de (4.2.10). Por (4.2.11) y el lema 4.3,
basta probar
\[
g_*\!\left(g^*(x)\left[(\rho i)^*c_1(\xi_{\mathbb P^1})
                       -g^*c_1(\xi_\Delta)\right]\right)=x.
\]
Por (4.2.4), $g_*g^*=0$, mientras que
\[
\begin{tikzcd}[column sep=huge]
\widetilde\Delta\arrow[r,"g\times\rho_i","\sim"']
  &\Delta\times\mathbb P^1,
\end{tikzcd}
\]
de modo que
\begin{equation}\tag{4.3.1}
g_*\!\left(g^*(x)(\rho i)^*c_1(\xi_{\mathbb P^1})\right)=x,
\end{equation}
y, por tanto,
\begin{equation}\tag{4.3.2}
g_*\!\left(g^*(x)c_1(\underline N_{\widetilde\Delta/\widetilde X})\right)
=-x.
\end{equation}

Calculemos ahora la estructura de álgebra sobre
\[
H^\bullet(X,\mathbb Z_\ell)\oplus
H^{\bullet-2}(\Delta,\mathbb Z_\ell(-1))
\]
transportada desde el producto cup mediante (4.2.2). Sean $a,b$
elementos de $H^\bullet(X,\mathbb Z_\ell)$ y $x,y$ elementos de
$H^{\bullet-2}(\Delta,\mathbb Z_\ell(-1))$. Comencemos con (4.2.3.1)
y escribamos
\[
(0\oplus x)\wedge(0\oplus y)=A\oplus B.
\]
Por (4.2.12) y (4.3.2),
\[
\begin{aligned}
A
&=f_*\!\left((i_*g^*x)(i_*g^*y)\right)\\
&=f_*i_*\!\left(g^*(xy)c_1(\underline N_{\widetilde\Delta/\widetilde X})\right)\\
&=j_*g_*\!\left(g^*(xy)c_1(\underline N_{\widetilde\Delta/\widetilde X})\right)
=-j_*(xy).
\end{aligned}
\]

\pageoriginale

Asimismo,
\[
\begin{aligned}
B
&=-g_*i^*\!\left((i_*g^*x)(i_*g^*y)\right)\\
&=-g_*i^*i_*\!\left(g^*(xy)c_1(\underline N_{\widetilde\Delta/\widetilde X})\right)\\
&=-g_*\!\left(g^*(xy)c_1(\underline N_{\widetilde\Delta/\widetilde X})^2\right).
\end{aligned}
\]
El lema 4.3 da
\[
\begin{aligned}
c_1(\underline N_{\widetilde\Delta/\widetilde X})^2
&=g^*\!\left(c_1(\xi_\Delta)^2\right)
 -2g^*c_1(\xi_\Delta)(\rho i)^*c_1(\xi_{\mathbb P^1}),
\end{aligned}
\]
porque $c_1(\xi_{\mathbb P^1})^2=0$. Por (4.2.4) y (4.3.1),
\[
B=0+2xy\,c_1(\xi_\Delta),
\]
lo que demuestra (4.2.3.1). La fórmula (4.2.3.2) expresa que $f^*$ es
un homomorfismo de álgebras.

Para demostrar (4.2.3.3), la fórmula de proyección muestra que
$\operatorname{Ker}f_*$ es estable por multiplicación por
$\operatorname{im}f^*$. Por tanto,
\[
(a\oplus0)\wedge(0\oplus y)=0\oplus C,
\]
donde
\begin{equation}\tag{4.3.3}
C=-g_*i^*\!\left(f^*(a)i_*g^*(y)\right).
\end{equation}
Usando (4.3.2),
\[
\begin{aligned}
C
&=-g_*\!\left(i^*f^*(a)g^*(y)
      c_1(\underline N_{\widetilde\Delta/\widetilde X})\right)\\
&=-g_*\!\left(g^*j^*(a)g^*(y)
      c_1(\underline N_{\widetilde\Delta/\widetilde X})\right)\\
&=j^*(a)y.
\end{aligned}
\]
Esto demuestra (4.2.3.3). Finalmente, (4.2.3.4) se deduce de
(4.2.3.3) por la conmutatividad graduada del producto cup.

\pageoriginale

\section*[5. Morfismo asociado a un pincel]{\texorpdfstring{5. Estudio cohomológico del morfismo
$\rho:\widetilde X\longrightarrow\mathbb P^1$ asociado a un pincel}{5. Estudio cohomológico del morfismo rho de X tilde a P1 asociado a un pincel}}

\subsection*{\texorpdfstring{5.1. El homomorfismo de restricción
$H^*(\widetilde X)\longrightarrow H^*(Y)$}{5.1. El homomorfismo de restricción en cohomología}}

Denotemos por $Y$ una de las fibras lisas de
$\rho:\widetilde X\to\mathbb P^1$, digamos
\[
Y=\rho^{-1}(t_0)=X\mathbin{\cdot}H_{t_0}.
\]
Se tiene el diagrama conmutativo
\begin{center}
\resizebox{\textwidth}{!}{%
\begin{tikzcd}[ampersand replacement=\&,column sep=large,row sep=huge]
\mathbb P^1
  \& {} \& \widetilde X\arrow[ll,"\rho"']\arrow[d,"f"]
  \& {} \& \widetilde\Delta
       \arrow[ll,hook',"i"]\arrow[d,"g"]
       \arrow[r,phantom,"\simeq"]
  \& \Delta\times\mathbb P^1 \\
\mathrm{Sp}(k)=\{t_0\}\arrow[u,hook]
  \& {} \& X
  \& {} \& \Delta
       \arrow[ll,hook',"j"]
       \arrow[ur,hook,"{\ell=\mathrm{id}\times\{t_0\}}"']
       \arrow[dlll,hook,"h"]
  \& {} \\
{} \& Y\arrow[ul]\arrow[uur,"k"]\arrow[ur,"m"']
  \& {} \& {} \& {} \& {}
\end{tikzcd}}
\end{center}

\pageoriginale

\subsection*{Proposición 5.1.1.}
\underline{Bajo las hipótesis de (3.1), designando $Y$ una fibra lisa:}

\smallskip
\noindent
\textup{(5.1.2)} \underline{El homomorfismo de restricción}
\[
k^*:H^\bullet(\widetilde X,\mathbb Z_\ell)
   \longrightarrow H^\bullet(Y,\mathbb Z_\ell)
\]
\underline{se expresa, mediante (4.2.2), por}
\[
H^\bullet(X,\mathbb Z_\ell)\oplus
H^{\bullet-2}(\Delta,\mathbb Z_\ell(-1))
\xrightarrow{\ m^*+\sisi{h^*}{h_*}\ }
H^\bullet(Y,\mathbb Z_\ell).
\]
\footnote{La impresión da $m^*+h^*$ en (5.1.2); el grado y la fórmula
de intersección transversal (5.1.5) exigen el morfismo de Gysin $h_*$.}

\noindent
\textup{(5.1.3)} \underline{La aplicación de Gysin}
\[
k_*:H^{\bullet-2}(Y,\mathbb Z_\ell(-1))
   \longrightarrow H^\bullet(\widetilde X,\mathbb Z_\ell)
\]
\underline{se expresa por}
\[
H^{\bullet-2}(Y,\mathbb Z_\ell(-1))
\xrightarrow{\ m_*\oplus(-h^*)\ }
H^\bullet(X,\mathbb Z_\ell)\oplus
H^{\bullet-2}(\Delta,\mathbb Z_\ell(-1)).
\]

\noindent\underline{Demostración.}
Los subesquemas lisos $\widetilde\Delta$ e $Y$ de $\widetilde X$ se
cortan transversalmente en $\Delta$; en otras palabras, el cuadrado de
esquemas lisos
\begin{equation}\tag{5.1.4}
\begin{tikzcd}[column sep=large,row sep=large]
\Delta\arrow[r,"\ell"]\arrow[d,hook,"h"]
  &\widetilde\Delta\arrow[d,hook,"i"]\\
Y\arrow[r,"k"']&\widetilde X
\end{tikzcd}
\end{equation}
es cartesiano. Por tanto, implica la fórmula [SGA 5, IV]
\begin{equation}\tag{5.1.5}
h_*\ell^*=k^*i_* ,\qquad \ell_*h^*=i^*k_*.
\end{equation}

\pageoriginale

El compuesto
\[
H^\bullet(X,\mathbb Z_\ell)\oplus
H^{\bullet-2}(\Delta,\mathbb Z_\ell(-1))
\xrightarrow{\ f^*+i_*g^*\ }
H^\bullet(\widetilde X,\mathbb Z_\ell)
\xrightarrow{\ k^*\ }
H^\bullet(Y,\mathbb Z_\ell)
\]
es la suma de $k^*f^*=m^*$ y
$k^*i_*g^*=h_*\ell^*g^*$, por (5.1.5). Como
$g\ell=\mathrm{id}_\Delta$, esto demuestra (5.1.2).

El compuesto
\[
H^{\bullet-2}(Y,\mathbb Z_\ell(-1))
\xrightarrow{\ k_*\ }
H^\bullet(\widetilde X,\mathbb Z_\ell)
\xrightarrow{\ f_*\oplus(-g_*i^*)\ }
H^\bullet(X,\mathbb Z_\ell)\oplus
H^{\bullet-2}(\Delta,\mathbb Z_\ell(-1))
\]
es la suma de
\[
f_*k_*=m_*
\qquad\text{y}\qquad
-g_*i^*k_*=-g_*\ell_*h^*=-h^*,
\]
lo que demuestra (5.1.3).

\subsection*{Corolario 5.1.6.}
\underline{Bajo las hipótesis de 5.1.1, para $i\leq n-1$, las imágenes
de}
\[
m^*:H^i(X,\mathbb Z_\ell)\longrightarrow H^i(Y,\mathbb Z_\ell)
\]
\underline{y de}
\[
k^*:H^i(\widetilde X,\mathbb Z_\ell)
   \longrightarrow H^i(Y,\mathbb Z_\ell)
\]
\underline{son iguales.}

\noindent\underline{Demostración.}
Por el teorema de Lefschetz débil para secciones hiperplanas
(SGA 5, VII, 7.1), para $i\leq n-1$ la aplicación
\begin{equation}\tag{5.1.7}
H^{i-2}(X,\mathbb Z_\ell(-1))
\xrightarrow{\ j^*\ }
H^{i-2}(\Delta,\mathbb Z_\ell(-1))
\end{equation}
es un isomorfismo. Así, por (5.1.2), la imagen de $k^*$ es la suma de
la imagen de $m^*$ y de la imagen del compuesto
\[
H^{i-2}(X,\mathbb Z_\ell(-1))
\xrightarrow{\ j^*\ }
H^{i-2}(\Delta,\mathbb Z_\ell(-1))
\xrightarrow{\ h_*\ }
H^i(Y,\mathbb Z_\ell).
\]

\pageoriginale

Pero la fórmula de proyección da
\[
\begin{aligned}
h_*j^*(x)
 &=h_*\bigl(h^*m^*(x)\bigr)
  =m^*(x)h_*(1)
  =\xi_Y\,m^*(x)\\
 &=m^*(x\xi_X),
 \qquad\text{pues }m^*(\xi_X)=\xi_Y.
\end{aligned}
\]
Por consiguiente, $\operatorname{Im}k^*\subset\operatorname{Im}m^*$,
y la inclusión inversa resulta directamente de (5.1.2).

\subsection*{5.2. El teorema de Lefschetz fuerte (LV)}

\subsubsection*{5.2.1.}
Sea $S$ proyectivo de dimensión $n$ sobre el cuerpo algebraicamente
cerrado $k$, y sea $S\hookrightarrow\mathbb P$ una inmersión
proyectiva. Denotemos por
\[
u\in H^2(S,\mathbb Z_\ell(1))
\]
la clase de cohomología de una sección hiperplana. Así, $u$ es la
imagen inversa por $S\hookrightarrow\mathbb P$ de la primera clase de
Chern de $\mathcal O_{\mathbb P}(1)$ en
$H^2(\mathbb P,\mathbb Z_\ell(1))$. Escribamos
\[
L_S:H^\bullet(S,\mathbb Q_\ell(j))
   \longrightarrow H^{\bullet+2}(S,\mathbb Q_\ell(j+1))
\]
para la multiplicación por $u$.

\subsubsection*{5.2.2.}
\underline{Se dice que $S$, provisto de $u$, satisface (LV), el teorema
de Lefschetz fuerte, si, para todo (o para un) entero $j$ y todo entero
$i>0$, las aplicaciones}
\[
L_S^i:H^{n-i}(S,\mathbb Q_\ell(j))
\longrightarrow H^{n+i}(S,\mathbb Q_\ell(i+j))
\]
\underline{son isomorfismos.}

\pageoriginale

\subsubsection*{5.2.2.1.}
Se conjetura que esta condición se cumple cuando $S$ es liso. Es
verdad cuando $k$ tiene característica cero [6]. El teorema de
especialización [SGA 4, XVI, 2.2] da entonces la misma conclusión
cuando $S$, junto con su inmersión proyectiva, se levanta a
característica cero. En particular, se cumple para una intersección
completa lisa. También se sabe para $S$ liso de dimensión a lo sumo
dos [2], y para una variedad abeliana, ya sea directamente
(Lieberman [4] y [2]), ya sea por el teorema de Mumford [5], según el
cual una variedad abeliana polarizada se levanta, junto con su
polarización, a característica cero.

Recordemos que, para una variedad proyectiva $S$ y $i\leq n$, la parte
primitiva de la cohomología se define por
\begin{equation}\tag{5.2.2.2}
\operatorname{Prim}^i(S,\mathbb Q_\ell(j))
=\operatorname{Ker}\!\left(
L^{n-i+1}:H^i(S,\mathbb Q_\ell(j))
\longrightarrow
H^{2n-i+2}\bigl(S,\mathbb Q_\ell(
  \sisi{j+2n-2i+2}{j+n-i+1})\bigr)
\right).
\end{equation}
\footnote{La potencia $L^{n-i+1}$ aumenta el giro de Tate en $n-i+1$; la
impresión da en el blanco el giro de Tate $j+2n-2i+2$.}
Si $S$ satisface (LV), para $i\leq n$ se tiene la descomposición
\begin{equation}\tag{5.2.2.3}
H^i(S,\mathbb Q_\ell(j))
\xleftarrow[\sim]{\ 1+L\ }
\operatorname{Prim}^i(S,\mathbb Q_\ell(j))
\oplus H^{i-2}(S,\mathbb Q_\ell(j-1)).
\end{equation}

\subsection*{Proposición 5.2.3.}
\underline{Sea $X$ liso y proyectivo y supongamos que satisface (LV).
Sea $\tau:Y\hookrightarrow X$ la inclusión de una sección hiperplana
lisa. Entonces}
\begin{equation}\tag{5.2.3.1}
\operatorname{Prim}^n(X,\mathbb Q_\ell(j))
=\operatorname{Ker}\!\left(
\tau^*:H^n(X,\mathbb Q_\ell(j))
\longrightarrow H^n(Y,\mathbb Q_\ell(j))
\right).
\end{equation}
\noindent\underline{Demostración.}
La fórmula de proyección para $\tau:Y\to X$ da el diagrama conmutativo

\pageoriginale

\begin{equation}\tag{5.2.3.2}
\begin{tikzcd}[column sep=huge,row sep=huge]
H^n(X,\mathbb Q_\ell(j))
  \arrow[rr,"L"]\arrow[dr,"\tau^*"']
  & {} &
H^{n+2}(X,\mathbb Q_\ell(j+1))
  \arrow[dl,"\tau_*"]\\
{} & H^n(Y,\mathbb Q_\ell(j)) & {}
\end{tikzcd}.
\end{equation}
Por dualidad y Lefschetz débil,
\[
\tau_*:H^n(Y,\mathbb Q_\ell(j))
\longrightarrow H^{n+2}(X,\mathbb Q_\ell(j+1))
\]
es la transpuesta de un isomorfismo y, por tanto, es un isomorfismo.
De aquí resulta la afirmación.

\subsection*{Lema 5.2.4.}
\underline{Sean $X$ liso y proyectivo y
$\tau:Y\hookrightarrow X$ la inclusión de una sección hiperplana lisa.
Entonces $X$ satisface (LV) si y solo si se cumplen las dos condiciones
siguientes:}

\smallskip
\noindent
\textup{5.2.4.1.} \underline{$Y$ satisface (LV).}

\smallskip
\noindent
\textup{5.2.4.2.} \underline{La restricción del producto cup sobre
$H^{n-1}(Y,\mathbb Q_\ell)$ al subespacio
$\operatorname{Im}H^{n-1}(X,\mathbb Q_\ell)$ es no degenerada.}

\noindent\underline{Demostración.}
Se tiene
\begin{equation}\tag{5.2.4.3}
\left\{
\begin{aligned}
L_X&=\tau_*\tau^*,\\[1ex]
L_Y&=\tau^*\tau_* \qquad\text{(por 4.2.11),}
\end{aligned}
\right.
\end{equation}

\pageoriginale

y, en consecuencia,
\begin{equation}\tag{5.2.4.4}
L_X^{i+1}=\tau_*L_Y^i\tau^*.
\end{equation}
En el diagrama conmutativo
\begin{equation}\tag{5.2.4.5}
\begin{tikzcd}[column sep=huge,row sep=huge]
H^{n-i-1}(X,\mathbb Q_\ell(j))
  \arrow[r,"L_X^{i+1}"]\arrow[d,"\tau^*"]
  &H^{n+i+1}(X,\mathbb Q_\ell(j+i+1))\\
H^{n-1-i}(Y,\mathbb Q_\ell(j))
  \arrow[r,"L_Y^i"]
  &H^{n-1+i}(Y,\mathbb Q_\ell(j+i))\arrow[u,"\tau_*"']
\end{tikzcd}
\end{equation}
$\tau^*$ es un isomorfismo para $i\geq1$ por Lefschetz débil, y
$\tau_*$ es un isomorfismo por dualidad. Por tanto, $X$ satisface
(LV) si y solo si $Y$ satisface (LV) y
\begin{equation}\tag{5.2.4.6}
L_X:H^{n-1}(X,\mathbb Q_\ell(j))
\longrightarrow H^{n+1}(X,\mathbb Q_\ell(j+1))
\end{equation}
es un isomorfismo.

Queda identificar (5.2.4.6) con 5.2.4.2. Sean
$a,b\in H^{n-1}(X,\mathbb Q_\ell(j))$. Por dualidad de Poincaré sobre
$X$, (5.2.4.6) equivale a la no degeneración de la forma bilineal
\begin{equation}\tag{5.2.4.7}
\begin{aligned}
H^{n-1}(X,\mathbb Q_\ell(j))\times
H^{n-1}(X,\mathbb Q_\ell(j))
&\longrightarrow H^{2n}\bigl(X,\mathbb Q_\ell(
  \sisi{2j+2}{2j+1})\bigr),\\
(a,b)&\longmapsto a\,L_X(b).
\end{aligned}
\end{equation}
Pero, por (5.2.4.3) y la fórmula de proyección,
\[
aL_X(b)=a\tau_*\tau^*(b)
=\tau_*\bigl(\tau^*(a)\tau^*(b)\bigr).
\]
Esto demuestra la equivalencia.

\pageoriginale

\subsection*{Corolario 5.2.5.}
\underline{Sea $X$ liso y proyectivo de dimensión $n$, y sea
$\tau:Y\hookrightarrow X$ la inclusión de una sección hiperplana lisa.
Supongamos que $X$ satisface (LV). Entonces, para $q\geq n$, la
aplicación}
\[
\tau^*:H^q(X,\mathbb Q_\ell(j))
\longrightarrow H^q(Y,\mathbb Q_\ell(j))
\]
\underline{es sobreyectiva.}

\noindent\underline{Demostración.}
Por (5.2.4.3), para $i\geq0$ se tiene
\[
L_Y^{i+1}=\tau^*L_X^i\tau_*,
\]
y, por tanto, el diagrama conmutativo
\begin{center}
\begin{tikzcd}[column sep=huge,row sep=huge]
H^{n-i}(X,\mathbb Q_\ell(j+1))
  \arrow[r,"L_X^i"]
  &H^{\sisi{n-i}{n+i}}(X,\mathbb Q_\ell(j+i+1))
    \arrow[d,"\tau^*"]\\
H^{n-i-2}(Y,\mathbb Q_\ell(j))
  \arrow[u,"\tau_*"']\arrow[r,"L_Y^{i+1}"]
  &H^{n+i}(Y,\mathbb Q_\ell(j+i+1)).
\end{tikzcd}
\end{center}
Esto da el resultado buscado.

\pageoriginale

\subsection*{Lema 5.2.6.}
\underline{Sea $f:F\to B$ un morfismo liso y propio, con $B$ conexo.
Sea $\mathcal O_F(1)$ un haz invertible relativamente amplio sobre
$F$, y dotemos cada fibra geométrica $F_{\bar b}$ del haz invertible
amplio inducido por $\mathcal O_F(1)$. Si una fibra geométrica
satisface (LV), todas las demás también lo satisfacen.}

\noindent\underline{Demostración.}
Sea $\xi_F\in H^2(F,\mathbb Z_\ell(1))$ la primera clase de Chern de
$\mathcal O_F(1)$. Esta clase y sus potencias
$\xi_F^i\in H^{\sisi{ii}{2i}}(F,\mathbb Z_\ell(i))$ inducen, mediante
el producto cup, los homomorfismos
\begin{equation}\tag{5.2.6.1}
\xi_F^i:R^{n-i}f_*\mathbb Q_\ell(j)
\longrightarrow R^{n+i}f_*\mathbb Q_\ell(j+i),
\end{equation}
donde $n$ es la dimensión relativa de las fibras. Por cambio de base
propio [SGA 4, XII, 5.2], todas las fibras geométricas satisfacen (LV)
si y solo si todas las aplicaciones (5.2.6.1) son isomorfismos.
Por el teorema de especialización [SGA 4, XVI, 2.2], los haces
$R^if_*\mathbb Q_\ell(j)$ son constantes torcidos. Por tanto, las
aplicaciones (5.2.6.1) son isomorfismos si y solo si lo son sobre un
punto geométrico dado.

\subsection*{Corolario 5.2.7.}
\underline{Sea $X$ liso y proyectivo. Si una sección hiperplana lisa
satisface (LV), toda sección hiperplana lisa lo satisface.}

\noindent\underline{Demostración.}
La familia universal de secciones hiperplanas lisas de $X$ satisface
las hipótesis del \sisi{lema 5.2.5}{lema 5.2.6}.

\pageoriginale

\subsection*{5.3. La condición (A)}

\subsubsection*{5.3.1.}
Sea
\[
\rho:\widetilde X\longrightarrow\mathbb P^1
\]
la ``fibración'' asociada a un pincel cuyo eje corta transversalmente
a $X$ y cuya fibra genérica se supone lisa.

Sea $U$ un abierto no vacío de $\mathbb P^1$ sobre el cual $\rho$ es
liso, y sea
\begin{equation}\tag{5.3.1.1}
\nu:U\hookrightarrow\mathbb P^1
\end{equation}
la inclusión. Como los haces $R^i\rho_*\mathbb Q_\ell(j)$ son
\underline{constantes torcidos sobre $U$}, los haces
\[
\nu_*\nu^*R^i\rho_*\mathbb Q_\ell(j)
\]
son esencialmente \underline{independientes} de la elección de $U$;
los homomorfismos canónicos de ``adjunción''
\begin{equation}\tag{5.3.2}
R^i\rho_*\mathbb Q_\ell(j)
\longrightarrow
\nu_*\nu^*R^i\rho_*\mathbb Q_\ell(j)
\end{equation}
también son independientes de la elección de $U$.

\subsubsection*{5.3.3.}
Obsérvese que, para un pincel \underline{arbitrario} de secciones
hiperplanas, el teorema de Lefschetz débil implica que, para
$i\leq n-2$, la inmersión
\[
\widetilde X\overset{(f\times\rho)}{\hookrightarrow}
X\times\mathbb P^1
\]
da un isomorfismo de haces sobre $\mathbb P^1$
\[
R^i\rho_*\mathbb Q_\ell(j)
\xleftarrow{\ \sim\ }
H^i(X,\mathbb Q_\ell(j))_{\mathbb P^1},
\qquad i\leq n-2.
\]

\pageoriginale

En otras palabras, $R^i\rho_*\mathbb Q_\ell(j)$ es constante para
$i\leq n-2$ y, con mayor razón, las aplicaciones (5.3.2) son
isomorfismos para $i\leq n-2$.

\subsubsection*{5.3.4.}
La teoría local de los ciclos evanescentes y, en particular, la
\underline{fórmula de Picard--Lefschetz} [SGA 7, XV, 3.4], muestra que,
\underline{para un pincel de Lefschetz, los homomorfismos (5.3.2) son
isomorfismos para $i\ne n$ y son sobreyectivos para $i=n$}.

\subsubsection*{5.3.5.}
Se dice que $\rho$, o el pincel del que procede,
\underline{satisface la condición (A)} si \underline{todos} los
homomorfismos (5.3.2) son isomorfismos; equivalentemente, si lo son
para todo $i$ con $j=0$.

Veremos más adelante (6.3, 6.4) que la condición (A) tiene una fuerte
tendencia a cumplirse.

\subsection*{Lema 5.4.}
\underline{Supongamos que una sección hiperplana lisa de $X$ (y por
tanto todas, por el corolario 5.2.7) satisface (LV), por ejemplo porque
$X$ satisface (LV), cf. el lema 5.2.4, y supongamos que $\rho$ satisface
la condición (A). Entonces toda fibra, incluidas las fibras singulares,
satisface (LV), y los haces $R^i\rho_*\mathbb Q_\ell(j)$ son constantes
para $i\ne n-1$.}

\noindent\underline{Demostración.}
Denotemos por
\[
L:R^i\rho_*\mathbb Q_\ell(j)
\longrightarrow R^{i+2}\rho_*\mathbb Q_\ell(j+1)
\]
la multiplicación por
$f^*c_1(\mathcal O_X(1))\in
H^2(\widetilde X,\mathbb Z_\ell(1))$; fibra por fibra, esta es la
operación considerada en (5.2.1). Por hipótesis, sobre $U$ se tiene

\pageoriginale

el isomorfismo
\begin{equation}\tag{5.4.1}
\nu^*(L^i):
\nu^*R^{n-1-i}\rho_*\mathbb Q_\ell(j)
\overset{\sim}{\longrightarrow}
\nu^*R^{n-1+i}\rho_*\mathbb Q_\ell(j+i)
\end{equation}
para todo $i>0$ y todo $j$. Aplicando el funtor $\nu_*$ se obtienen los
\underline{isomorfismos}
\begin{equation}\tag{5.4.2}
(\nu_*\nu^*)L^i:
\nu_*\nu^*R^{n-1-i}\rho_*\mathbb Q_\ell(j)
\overset{\sim}{\longrightarrow}
\nu_*\nu^*R^{n-1+i}\rho_*\mathbb Q_\ell(j+i).
\end{equation}
En virtud de (A), estos se identifican con los morfismos
\begin{equation}\tag{5.4.3}
L^i:R^{n-1-i}\sisi{\rho^*}{\rho_*}\mathbb Q_\ell(j)
\overset{\sim}{\longrightarrow}
R^{n-1+i}\rho_*\mathbb Q_\ell(j+i).
\end{equation}
La conclusión resulta de 5.3.3.

\subsubsection*{5.4.4.}
Sea $Y$ una fibra lisa de
$\rho:\widetilde X\to\mathbb P^1$. Por Lefschetz débil, la aplicación
de restricción
\[
H^{n-1}(X,\mathbb Q_\ell(j))
\overset{m^*}{\hookrightarrow}
H^{n-1}(Y,\mathbb Q_\ell(j))
\]
es inyectiva. Vimos en 5.2.4 que, si $X$ satisface (LV), el producto
cup sobre $H^{n-1}(Y,\mathbb Q_\ell(j))$ sigue siendo no degenerado al
restringirlo a $m^*H^{n-1}(X,\mathbb Q_\ell(j))$. El complemento
ortogonal de $H^{n-1}(X,\mathbb Q_\ell(j))$ en
$H^{n-1}(Y,\mathbb Q_\ell(j))$ es, por tanto, un subespacio
complementario; siguiendo la terminología clásica, se llama
\underline{espacio de cohomología evanescente de $Y$}.

Sea $\nu:U\hookrightarrow\mathbb P^1$ un abierto no vacío sobre el que
$\rho$ es liso. El subhaz correspondiente de
$\nu_*\nu^*R^{n-1}\rho_*\mathbb Q_\ell(j)$, el
\underline{haz de cohomología evanescente}
$E^{n-1}\rho_*\mathbb Q_\ell(j)$ sobre $\mathbb P^1$, se define por

\pageoriginale

\begin{equation}\tag{5.4.4.1}
E^{n-1}\rho_*\mathbb Q_\ell(j)
=\nu_*\left\{
\begin{array}{l}
\text{el complemento ortogonal en }
\nu^*R^{n-1}\rho_*\mathbb Q_\ell(j)\text{ del}\\[1ex]
\text{subhaz constante }H^{n-1}(X,\mathbb Q_\ell(j))_U
\end{array}
\right\}.
\end{equation}
Es claro que este haz es independiente de la elección de $U$.

\subsection*{Lema 5.5.}
\underline{Supongamos que $X$ satisface (LV) y que las aplicaciones
(5.3.2) son isomorfismos para $i=n-1$, por ejemplo porque $\rho$
procede de un pincel de Lefschetz, cf. 5.3.4. Entonces se tiene una
descomposición en suma directa}
\begin{equation}\tag{5.5.1}
R^{n-1}\rho_*\mathbb Q_\ell(j)
\simeq
H^{n-1}(X,\mathbb Q_\ell(j))_{\mathbb P^1}
\oplus E^{n-1}\rho_*\mathbb Q_\ell(j).
\end{equation}
\noindent\underline{Demostración.}
Por (LV), sobre $U$ se tiene una descomposición
\begin{equation}\tag{5.5.2}
\nu^*R^{n-1}\rho_*\mathbb Q_\ell(j)
\simeq
H^{n-1}(X,\mathbb Q_\ell(j))_U
\oplus \nu^*E^{n-1}\rho_*\mathbb Q_\ell(j),
\end{equation}
a la que se aplica $\nu_*$.

\subsection*{Teorema 5.6.}
\underline{Sea $X$ una variedad proyectiva, lisa e irreducible que
satisface (LV) (5.2.2). Sea $\{H_\tau\}$ un pincel de hiperplanos cuyo
miembro genérico y cuyo eje cortan transversalmente a $X$, en $X_\eta$
y en $\Delta$, respectivamente. Supongamos que la proyección
$\rho:\widetilde X\to\mathbb P^1$ satisface la
\sisi{condición (5.3.2)}{condición (A)}. Entonces:}
\footnote{La impresión dice ``condición (5.3.2)'', pero (5.3.2) es la
familia de aplicaciones de adjunción; la condición (A) se define a
partir de ellas en 5.3.5.}

\smallskip
\noindent
\textup{(5.6.1)}\quad
\underline{la sucesión espectral de Leray}
\[
E_2^{p,q}=H^p(\mathbb P^1,R^q\rho_*\mathbb Q_\ell)
\Longrightarrow H^{\sisi{m+q}{p+q}}(\widetilde X,\mathbb Q_\ell)
\]
\underline{degenera.}
\footnote{La convergencia impresa lleva $H^{m+q}$; el bigrado y el
enunciado paralelo del teorema 5.6.8 exigen $H^{p+q}$.}

\noindent
\textup{(Nota.)}\quad
\underline{Esto sigue siendo válido si, en vez de suponer que $X$
satisface (LV), solo se supone que lo satisface una sección hiperplana
lisa $Y$.}

\pageoriginale

\subsubsection*{5.6.2.}
\underline{Denotemos por $\eta$ el punto genérico de $\mathbb P^1$,
por $\bar\eta$ el punto genérico geométrico, por $X_{\bar\eta}$ la fibra
genérica geométrica de $\rho$, y por $\pi$ el grupo de Galois de
$k(\bar\eta)/k(\eta)$.}

\underline{Para $q\ne n-1$ se tienen isomorfismos canónicos}
\begin{align}
E_2^{\sisi{o}{0},q}
 &=H^0(\mathbb P^1,R^q\rho_*\mathbb Q_\ell)
 \simeq H^q(X_{\bar\eta},\mathbb Q_\ell),
 \tag{5.6.2.1}\\
E_2^{2,q}
 &=H^2(\mathbb P^1,R^q\rho_*\mathbb Q_\ell)
 \simeq H^q(X_{\bar\eta},\mathbb Q_\ell(-1)),
 \tag{5.6.2.2}\\
E_2^{p,q}
 &=H^p(\mathbb P^1,R^q\rho_*\mathbb Q_\ell)=0,
 \qquad p\ne0,2.
 \tag{5.6.2.3}
\end{align}
\underline{Para $q=n-1$ se tiene}
\begin{align}
E_2^{0,n-1}
 &=H^0(\mathbb P^1,R^{n-1}\rho_*\mathbb Q_\ell)
 \simeq H^{n-1}(X_{\bar\eta},\mathbb Q_\ell)^\pi
 \simeq H^{n-1}(X,\mathbb Q_\ell),
 \tag{5.6.2.4}\\
E_2^{2,n-1}
 &=H^2(\mathbb P^1,R^{n-1}\rho_*\mathbb Q_\ell)
 \simeq H^{n-1}(X_{\bar\eta},\mathbb Q_\ell(-1))^\pi
 \simeq H^{n-1}(X,\mathbb Q_\ell(-1)).
 \tag{5.6.2.5}
\end{align}
\underline{(El valor restante $E_2^{1,n-1}$ se hallará en 5.7.)}

\noindent\underline{Demostración.}
Por (5.4.3), se aplica el teorema 1.1 y se obtiene (5.6.1). Las tres
primeras fórmulas resultan de que $R^q\rho_*\mathbb Q_\ell$ es constante
para $q\ne n-1$ (lema 5.4); para (5.6.2.3) se usa además que
$\mathbb P^1$ es simplemente conexa.

\pageoriginale

El primer isomorfismo de (5.6.2.4) resulta de
\[
R^{n-1}\rho_*\mathbb Q_\ell
\overset{\sim}{\longrightarrow}
\nu_*\nu^*R^{n-1}\rho_*\mathbb Q_\ell.
\]
En efecto,
\[
H^0(\mathbb P^1,R^{n-1}\rho_*\mathbb Q_\ell)
\simeq H^0(U,\nu^*R^{n-1}\rho_*\mathbb Q_\ell),
\]
y $\nu^*R^{n-1}\rho_*\mathbb Q_\ell$ es un haz constante torcido sobre
$U$, por lo que corresponde al $\pi$-módulo
$H^{n-1}(X_{\bar\eta},\mathbb Q_\ell)$. El segundo isomorfismo de
(5.6.2.4) resulta de la degeneración junto con el corolario 5.1.6.

Pasemos a (5.6.2.5). Sea $F^\bullet$ la filtración sobre
$H^\bullet(\widetilde X,\mathbb Q_\ell)$ definida por la sucesión
espectral. Por degeneración,
$F^2H^{n+1}(\widetilde X,\mathbb Q_\ell)=E_2^{2,n-1}$ y
$\operatorname{gr}^1H^{n+1}(\widetilde X,\mathbb Q_\ell)=E_2^{1,n}=0$.
Por tanto,
\begin{align}
E_2^{2,n-1}
 &=F^1H^{n+1}(\widetilde X,\mathbb Q_\ell) \notag\\
 &=\operatorname{Ker}\!\left(
 H^{n+1}(\widetilde X,\mathbb Q_\ell)
 \longrightarrow H^0(\mathbb P^1,R^{n+1}\rho_*\mathbb Q_\ell)
 \right) \notag\\
 &\overset{(A)}{=}
 \operatorname{Ker}\!\left(
 k^*:H^{n+1}(\widetilde X,\mathbb Q_\ell)
 \longrightarrow H^{n+1}(X_{\bar\eta},\mathbb Q_\ell)
 \right).
\tag{5.6.3}
\end{align}
Por (5.1.2), $k^*$ se escribe
\begin{equation}\tag{5.6.4}
H^{n+1}(X,\mathbb Q_\ell)
\oplus H^{n-1}(\Delta,\mathbb Q_\ell(-1))
\xrightarrow{\ m^*+h_*\ }
H^{n+1}(X_{\bar\eta},\mathbb Q_\ell).
\end{equation}
Usando (5.2.2.3), se convierte en
\begin{equation}\tag{5.6.5}
\begin{split}
&\operatorname{Prim}^{n-1}(X,\mathbb Q_\ell(-1)\sisi{}{)}
\oplus H^{n-3}(X,\mathbb Q_\ell(-2))
\oplus H^{n-1}(\Delta,\mathbb Q_\ell(-1))\\
&\hspace{35mm}\xrightarrow{\ m^*L+m^*L^2\sisi{}{+}h_*\ }
H^{n+1}(X_{\bar\eta},\mathbb Q_\ell\sisi{}{)}.
\end{split}
\end{equation}
\footnote{La impresión de (5.6.5) omite varios paréntesis finales y el
signo más ante la aplicación $h_*$ del tercer sumando. Los tres
sumandos fuente y (5.6.4) fuerzan la aplicación de tres términos.}

La aplicación
\[
H^{n-3}(X,\mathbb Q_\ell(-2))
\xrightarrow{\ m^*L^2\ }
H^{n+1}(X_{\bar\eta},\mathbb Q_\ell)
\]
es un isomorfismo por (LV), ya que
$m^*L_X^2=L_{X_{\bar\eta}}^2m^*$. De ahí que el núcleo de
\sisi{(5.6.8)}{(5.6.5)} se identifique, mediante
$\operatorname{pr}_1\oplus\operatorname{pr}_3$, con

\pageoriginale

\[
\operatorname{Prim}^{n-1}(X,\mathbb Q_\ell(-1))
\oplus H^{n-1}(\Delta,\mathbb Q_\ell(-1)).
\]
Por (LV) y Lefschetz débil,
\[
H^{n-1}(\Delta,\mathbb Q_\ell(-1))
\xleftarrow[\sim]{\,L\,}
H^{n-3}(\Delta,\mathbb Q_\ell(-2))
\xleftarrow[\sim]{\ }
H^{n-3}(X,\mathbb Q_\ell(-2)).
\]
Junto con (5.2.2.3), esto da
\begin{equation}\tag{5.6.6}
E_2^{2,n-1}\simeq H^{n-1}(X,\mathbb Q_\ell(-1)).
\end{equation}
El otro isomorfismo de (5.6.2.5) se deduce torciendo (5.6.2.4).

\subsubsection*{Observación 5.6.7.}
La afirmación (5.6.1) puede demostrarse bajo hipótesis bastante más
débiles, como me explicó Deligne:

\subsection*{Teorema 5.6.8.}
\underline{Sea $X$ liso de dimensión $n$ sobre un cuerpo
algebraicamente cerrado $k$, y sea}
\[
\rho:X\longrightarrow\mathbb P^1_k
\]
\underline{un morfismo proyectivo cuya fibra genérica es lisa.
Supongamos que:}

\smallskip
\noindent
\textup{(5.6.8.1)}\quad
\underline{la fibra genérica geométrica, dotada de la restricción del
haz invertible relativamente amplio $\mathcal O_X(1)$, satisface (LV);}

\smallskip
\noindent
\textup{(5.6.8.2)}\quad
\underline{las aplicaciones (5.3.2) son isomorfismos para $i\ne n$ y
son sobreyectivas para $i=n$ (condición que se satisface, en
particular, si $\rho$ procede de un pincel de Lefschetz, cf. 5.3.4).
Entonces la sucesión espectral de Leray}
\[
E_2^{p,q}=H^p(\mathbb P^1,R^q\rho_*\mathbb Q_\ell)
\Longrightarrow H^{p+q}(X,\mathbb Q_\ell)
\]
\underline{degenera.}

\smallskip
\noindent
\underline{Observación.}
El mismo teorema, con la misma demostración, vale sustituyendo
$\mathbb P^1_k$ por una curva proyectiva, lisa e irreducible cualquiera.

\pageoriginale

\noindent\underline{Demostración.}
Para la demostración identificaremos el haz $\mathbb Q_\ell(1)$ con
$\mathbb Q_\ell$, lo que es posible porque $k$ es algebraicamente
cerrado. Podemos suponer que $X$ es conexo y que $\rho$ es
\underline{sobreyectivo}; por tanto, es plano, con fibras de dimensión
$n-1$. Sea
\[
\xi_X\in H^2(X,\mathbb Q_\ell)
\]
la primera clase de Chern de $\mathcal O_X(1)$, y sea
\[
\nu:U\hookrightarrow\mathbb P^1
\]
la inclusión de un abierto de lisura de $\rho$. Por (5.6.8.1) y el
lema 5.2.6, las aplicaciones de producto cup
\[
\xi_X^i:
\nu^*R^{n-1-i}\rho_*\mathbb Q_\ell
\longrightarrow
\nu^*R^{n-1+i}\rho_*\mathbb Q_\ell
\]
son isomorfismos para $i\geq1$. Aplicando $\nu_*$ y usando (5.6.8.2)
se obtienen, para $i\geq2$, los isomorfismos
\begin{equation}\tag{5.6.8.3}
\xi_X^i:R^{n-1-i}\rho_*\mathbb Q_\ell
\overset{\sim}{\longrightarrow}
R^{n-1+i}\rho_*\mathbb Q_\ell,
\end{equation}
y, para $i=1$, el isomorfismo
\begin{equation}\tag{5.6.8.4}
\xi_X:R^{n-2}\rho_*\mathbb Q_\ell
\overset{\sim}{\longrightarrow}
\nu_*\nu^*R^n\rho_*\mathbb Q_\ell.
\end{equation}
Este último se factoriza como
\begin{equation}\tag{5.6.8.5}
R^{n-2}\rho_*\mathbb Q_\ell
\xrightarrow{\ \xi_X\ }
R^n\rho_*\mathbb Q_\ell
\xrightarrow{\ \text{adjunción}\ }
\nu_*\nu^*R^n\rho_*\mathbb Q_\ell.
\end{equation}

\pageoriginale

Por consiguiente, se tiene una descomposición canónica
\begin{equation}\tag{5.6.8.6}
R^n\rho_*\mathbb Q_\ell
\xleftarrow[\sim]{\ \xi_X+\operatorname{incl}\ }
R^{n-2}\rho_*\mathbb Q_\ell\oplus
\operatorname{Local}^n\mathbb Q_\ell,
\end{equation}
donde
\begin{equation}\tag{5.6.8.7}
\operatorname{Local}^n\mathbb Q_\ell
=\operatorname{Ker}\!\left(
R^n\rho_*\mathbb Q_\ell
\longrightarrow\nu_*\nu^*R^n\rho_*\mathbb Q_\ell
\right).
\end{equation}
El haz $\operatorname{Local}^n\mathbb Q_\ell$ está concentrado en los
puntos de $\mathbb P^1$ donde la fibra de $\rho$ no es lisa.
Combinando este hecho con (5.6.8.3), se obtiene
\begin{equation}\tag{5.6.8.8}
\xi_X^{\sisi{1}{i+1}}:
H^j(\mathbb P^1,R^{n-2-i}\rho_*\mathbb Q_\ell)
\overset{\sim}{\longrightarrow}
H^j(\mathbb P^1,R^{\sisi{n-1+i}{n+i}}\rho_*\mathbb Q_\ell),
\qquad \sisi{i\geq1,\ j\geq1}{i\geq0,\ j\geq1}.
\end{equation}
\footnote{La impresión de (5.6.8.8) tiene un desplazamiento impar de
grado, etiqueta la aplicación solo como $\xi_X$ y comienza en $i=1$.
La forma corregida viene impuesta por (5.6.8.3), (5.6.8.6), la flecha
inferior del cuadrado siguiente y la especialización posterior
$E_2^{1,n-2}\simeq E_2^{1,n}$.}

Demostremos ahora la degeneración. Como $E_2^{p,q}=0$ salvo para
$p=0,1,2$, basta probar que todos los diferenciales
\[
d_2:H^0(\mathbb P^1,R^i\rho_*\mathbb Q_\ell)
\longrightarrow H^2(\mathbb P^1,R^{i-1}\rho_*\mathbb Q_\ell)
\]
son nulos. Para $q\leq n-1$, pongamos
\begin{equation}\tag{5.6.8.9}
\operatorname{Prim}(E_2^{p,q})
=\operatorname{Ker}\!\left(
\xi_X^{n-q}:E_2^{p,q}\longrightarrow E_2^{p,2n-q}
\right).
\end{equation}
Las fórmulas (5.6.8.3) y (5.6.8.6), y en particular la inyectividad de
\[
\xi_X:H^p(\mathbb P^1,R^{n-2}\rho_*\mathbb Q_\ell)
\longrightarrow H^p(\mathbb P^1,R^n\rho_*\mathbb Q_\ell),
\]
dan, para $q\leq n-1$, la descomposición
\begin{equation}\tag{5.6.8.10}
E_2^{p,q}
\xleftarrow[\sim]{\ 1+\xi_X+\xi_X^2+\cdots\ }
\operatorname{Prim}(E_2^{p,q})
\oplus\operatorname{Prim}(E_2^{p,q-2})
\oplus\operatorname{Prim}(E_2^{p,q-4})
\oplus\cdots.
\end{equation}

\pageoriginale

Para probar que
\[
d_2^{p,q}=0\qquad(q\leq n-1),
\]
consideremos el diagrama conmutativo
\[
\begin{tikzcd}[column sep=huge,row sep=huge]
\operatorname{Prim}(E_2^{p,q})
  \arrow[r,"\xi_X^{n-q}=0"]
  \arrow[d,"d_2^{p,q}"]
&
E_2^{p,2n-q}
  \arrow[d,"d_2^{p,2n-q}"]
\\
E_2^{p+2,q-1}
  \arrow[r,"\xi_X^{n-1-(q-1)}","\sim"']
&
E_2^{p+2,2n-q-1}.
\end{tikzcd}
\]
La flecha inferior es un isomorfismo por (5.6.8.8). El diagrama
muestra que $d_2^{p,q}$ anula $\operatorname{Prim}(E_2^{p,q})$ para
$q\leq n-1$, y (5.6.8.10) muestra entonces que $d_2$ anula
$E_2^{p,q}$ para $q\leq n-1$. Por (5.6.8.3), también lo anula para
$q\geq n+1$. Solo queda demostrar que
\[
d_2^{0,n}:E_2^{0,n}\longrightarrow E_2^{2,n-1}
\]
es nulo, o equivalentemente que
\begin{equation}\tag{5.6.8.12}
\dim E_3^{2,n-1}=\dim E_2^{2,n-1}.
\end{equation}
Ya hemos probado
\begin{equation}\tag{5.6.8.13}
\dim E_2^{p,q}=\dim E_3^{p,q}
\qquad\text{para }(p,q)\ne(0,n),(2,n-1),
\end{equation}
y, por razones de grado,
\begin{equation}\tag{5.6.8.14}
E_3^{p,q}=E_\infty^{p,q}
\qquad\text{para todo par }(p,q).
\end{equation}

\pageoriginale

De ello resulta
\[
b_i(X)=\dim H^i(X,\mathbb Q_\ell)
=\sum_{p+q=i}\dim E_3^{p,q}.
\]
Teniendo en cuenta (5.6.8.13), se obtiene
\begin{equation}\tag{5.6.8.15}
b_{n-1}(X)
=\dim E_2^{\sisi{c}{0},n-1}
 +\dim E_2^{1,n-2}
 +\sisi{\dim_2^{2,n-3}}{\dim E_2^{2,n-3}}
\end{equation}
y
\begin{equation}\tag{5.6.8.16}
b_{n+1}(X)
=\dim E_2^{0,n+1}
 +\dim E_2^{1,n}
 +\dim E_3^{2,n-1}.
\end{equation}
La dualidad de Poincaré da
\begin{equation}\tag{5.6.8.17}
b_{n-1}(X)=b_{n+1}(X).
\end{equation}
Por (5.6.8.8),
\[
\dim E_2^{1,n}=\dim E_2^{1,n-2}.
\]
Así, (5.6.8.17) muestra que basta conocer
\begin{align}
\dim E_2^{2,n-3}&=\dim E_2^{0,n+1},
\tag{5.6.8.18}\\
\dim E_2^{0,n-1}&=\dim E_2^{2,n-1}.
\tag{5.6.8.19}
\end{align}
Estas fórmulas resultan de (5.6.8.2), de la dualidad de Poincaré para
las fibras lisas de $\rho$ y del lema siguiente, aplicado a
$\mathbb P^1$ y a los haces $R^{n-3}\rho_*\mathbb Q_\ell$ y
$R^{n-1}\rho_*\mathbb Q_\ell$.

\pageoriginale

\subsection*{Lema 5.6.9.}
\underline{Sea $D$ una curva lisa, conexa y propia sobre un cuerpo
algebraicamente cerrado $k$. Sea $i:\eta\hookrightarrow D$ la inclusión
del punto genérico, y sea $F$ un $\mathbb Q_\ell$-haz constructible
sobre $D$ tal que
$F\overset{\sim}{\longrightarrow}i_*i^*F$. Entonces}
\[
H^2(D,F)\simeq H^0(D,\check F(-1))^\vee,
\]
\underline{donde $\check F\simeq i_*((i^*F)^\vee)$.}

\subsection*{Corolario 5.6.10.}
\textup{(Utiliza Picard--Lefschetz.)}
\underline{Sea $\rho:\widetilde X\to\mathbb P^1$ la ``fibración''
procedente de un pincel de Lefschetz, y supongamos que las fibras lisas
satisfacen (LV). Entonces}
\begin{equation}\tag{5.6.10.1}
H^i(X_{\bar\eta},\mathbb Q_\ell)^\pi
=\operatorname{Im}\!\left(
H^i(X,\mathbb Q_\ell)\longrightarrow
H^i(X_{\bar\eta},\mathbb Q_\ell)
\right).
\end{equation}
\noindent\underline{Demostración.}
Por 5.3.4, que utiliza la fórmula de Picard--Lefschetz, el teorema
5.6.8 se aplica y da
\[
H^i(X_{\bar\eta},\mathbb Q_\ell)^\pi
=\operatorname{Im}\!\left(
H^i(\widetilde X,\mathbb Q_\ell)\longrightarrow
H^i(X_{\bar\eta},\mathbb Q_\ell)
\right).
\]
Queda comprobar que $H^i(X,\mathbb Q_\ell)$ y
$H^i(\widetilde X,\mathbb Q_\ell)$ tienen la misma imagen en
$H^i(X_{\bar\eta},\mathbb Q_\ell)$. Esto vale para $i\leq n-1$ por el
corolario 5.1.6, y para $i\geq n$ por el corolario 5.2.5.

Pasemos ahora al cálculo más importante.

\subsection*{Teorema 5.7.}
\underline{Bajo las hipótesis de 5.6, pongamos}
\[
\begin{aligned}
E^{n-2}(\Delta,\mathbb Q_\ell(j))
&=\underline{\text{el complemento ortogonal en }}
H^{n-2}(\Delta,\mathbb Q_\ell(j))\\
&\qquad\underline{\text{de la imagen de }}
H^{n-2}(X,\mathbb Q_\ell(j)).
\end{aligned}
\]
\underline{Esta es la parte ``evanescente'' de la cohomología de
$\Delta$, cf. 5.4.4. Entonces se tiene una descomposición en suma
directa} (cf. 4.2)
\begin{equation}\tag{5.7.1}
E_2^{1,n-1}
=H^1(\mathbb P^1,R^{n-1}\rho_*\mathbb Q_\ell)
\xleftarrow[\;f^*+i_*g^*\;]{\sim}
\operatorname{Prim}^n(X,\mathbb Q_\ell)
\oplus E^{n-2}(\Delta,\mathbb Q_\ell(-1)).
\end{equation}

\pageoriginale

\noindent\underline{Demostración.}
Por degeneración,
\[
E_2^{1,n-1}\simeq
F^1H^n(\widetilde X,\mathbb Q_\ell)
\big/F^2H^n(\widetilde X,\mathbb Q_\ell).
\]
Por (A),
\begin{equation}\tag{5.7.2}
F^1H^n(\widetilde X,\mathbb Q_\ell)
=\ker\!\left(
k^*:H^n(\widetilde X,\mathbb Q_\ell)\longrightarrow
H^n(X_{\bar\eta},\mathbb Q_\ell)
\right),
\end{equation}
lo que, por (5.1.2), se identifica con
\[
\ker\!\left(
H^n(X,\mathbb Q_\ell)
\oplus H^{n-2}(\Delta,\mathbb Q_\ell(-1))
\xrightarrow{\;m^*+\sisi{h^*}{h_*}\;}
H^n(X_{\bar\eta},\mathbb Q_\ell)
\right).
\]
Bajo (LV), $k^*$ también se escribe
\begin{equation}\tag{5.7.3}
\operatorname{Prim}^n(X,\mathbb Q_\ell)
\oplus H^{n-2}(X,\mathbb Q_\ell(-1))
\oplus H^{n-2}(\Delta,\mathbb Q_\ell(-1))
\xrightarrow{\;m^*+m^*L+h_*\;}
H^n(X_{\bar\eta},\mathbb Q_\ell),
\end{equation}
y la aplicación central
\begin{equation}\tag{5.7.4}
H^{n-2}(X,\mathbb Q_\ell(-1))
\xrightarrow{\;m^*L\;}
H^n(X_{\bar\eta},\mathbb Q_\ell)
\end{equation}
es un isomorfismo (cf. 5.2.3 y (5.2.2.3)). Denotando por
\begin{equation}\tag{5.7.5}
\wp:H^n(X,\mathbb Q_\ell)\longrightarrow
\operatorname{Prim}^n(X,\mathbb Q_\ell)
\end{equation}
la proyección sobre la parte primitiva, obtenemos
\begin{equation}\tag{5.7.6}
F^1H^n(\widetilde X,\mathbb Q_\ell)
\xrightarrow[\;\wp f_*\oplus(-g_*\sisi{2^*}{i^*})\;]{\sim}
\operatorname{Prim}^n(X,\mathbb Q_\ell)
\oplus H^{n-2}(\Delta,\mathbb Q_\ell(-1)).
\end{equation}
\footnote{La etiqueta impresa de (5.7.6) lleva $2^*$ tras $g_*$.
La descomposición inversa de la explosión (4.2.2), utilizada aquí,
exige $i^*$.}
Por (5.6.2.2) y degeneración,
\begin{equation}\tag{5.7.7}
F^2H^n(\widetilde X,\mathbb Q_\ell)
\simeq E_2^{2,n-2}
\simeq H^{n-2}(X_{\bar\eta},\mathbb Q_\ell(-1)).
\end{equation}

\pageoriginale

Por otra parte, la inclusión de
$F^2H^{\sisi{N}{n}}(\widetilde X,\mathbb Q_\ell)
=H^{n-2}(X_{\bar\eta},\mathbb Q_\ell(-1))$
en $H^n(\widetilde X,\mathbb Q_\ell)$ es la aplicación de Gysin $k_*$.
En rigor, la aplicación de Gysin solo está definida como una aplicación
\[
\sisi{H^{i-2}(X_{\bar\eta},\mathbb Q_\ell)}
      {H^{i-2}(X_{\bar\eta},\mathbb Q_\ell(-1))}
\longrightarrow
H^i(X\times_k k(\bar\eta),\mathbb Q_\ell),
\]
pero, en virtud del isomorfismo canónico (SGA 4, XII, 5.4)
\[
H^i(X,\mathbb Q_\ell)\xrightarrow{\sim}
H^i(X\otimes_kK,\mathbb Q_\ell)
\]
para todo cuerpo algebraicamente cerrado $K/k$, esperamos que el lector
nos permita este abuso.

Por (5.1.3), la aplicación de Gysin se expresa como
\begin{equation}\tag{5.7.8}
H^{n-2}(X_{\bar\eta},\mathbb Q_\ell(-1))
\xrightarrow{\;m_*\oplus(-h^*)\;}
H^n(X,\mathbb Q_\ell)
\oplus H^{n-2}(\Delta,\mathbb Q_\ell(-1)).
\end{equation}
Por Lefschetz débil,
\[
m_*H^{n-2}(X_{\bar\eta},\mathbb Q_\ell(-1))
=m_*m^*H^{n-2}(X,\mathbb Q_\ell(-1))
=L_XH^{n-2}(X,\mathbb Q_\ell(-1))
\]
es el complemento ortogonal de la parte primitiva. Por tanto, se tiene
el diagrama conmutativo
\begin{equation}\tag{5.7.9}
\begin{tikzcd}[column sep=huge,row sep=huge]
F^2H^n(\widetilde X,\mathbb Q_\ell)
  \arrow[r,hook]
& F^1H^n(\widetilde X,\mathbb Q_\ell)
  \arrow[d,"\wp f_*\oplus(-g_*i^*)","\wr"']
\\
H^{n-2}(X_{\bar\eta},\mathbb Q_\ell(-1))
  \arrow[u,"k_*"',"\wr"]
  \arrow[r,"0\oplus h^*"]
& \operatorname{Prim}^n(X,\mathbb Q_\ell)
  \oplus H^{n-2}(\Delta,\mathbb Q_\ell(-1)).
\end{tikzcd}
\end{equation}
Esto demuestra (5.7.1).

\pageoriginale

\subsection*{Observación 5.7.10.}
En 5.7, la hipótesis (A) puede sustituirse por la hipótesis más débil
de que, para $q=n$, las aplicaciones
\begin{equation}\tag{5.7.11}
H^q(\widetilde X,\mathbb Q_\ell)\longrightarrow
H^0(\mathbb P^1,R^q\rho_*\mathbb Q_\ell)
\end{equation}
y
\begin{equation}\tag{5.7.12}
H^q(\widetilde X,\mathbb Q_\ell)\longrightarrow
H^q(X_{\bar\eta},\mathbb Q_\ell)
\end{equation}
tengan el mismo núcleo, como muestra la demostración. Sin esta
condición, la conclusión puede fallar; cf. 5.9.

\section*{5.8. El homomorfismo de Griffiths}

\subsection*{5.8.1.}
Sea $f:T\to S$ un morfismo de esquemas. La sucesión espectral de Leray
define una filtración $(F^i)$ sobre los grupos
$H^\bullet(T,\mathbb Q_\ell(j))$. Se define
\begin{equation}\tag{5.8.1.1}
\operatorname{Prim}^\bullet(T/S,\mathbb Q_\ell(j))
=F^1H^\bullet(T,\mathbb Q_\ell(j)).
\end{equation}
Como la sucesión espectral de Leray está situada en el primer cuadrante,
se tiene
\begin{equation}\tag{5.8.1.2}
E_\infty^{0,\bullet}\lhook\joinrel\longrightarrow E_2^{0,\bullet},
\end{equation}
y, por tanto,
\begin{equation}\tag{5.8.1.3}
\operatorname{Prim}^\bullet(T/S,\mathbb Q_\ell(j))
=\ker\!\left(
H^\bullet(T,\mathbb Q_\ell(j))
\longrightarrow H^0(S,R^\bullet f_*\mathbb Q_\ell(j))
\right).
\end{equation}
Por la misma razón,
\begin{equation}\tag{5.8.1.4}
E_\infty^{1,\bullet}\lhook\joinrel\longrightarrow E_2^{1,\bullet}.
\end{equation}

\pageoriginale

El compuesto
\begin{equation}\tag{5.8.1.5}
\begin{tikzcd}[column sep=large,row sep=large]
\operatorname{Prim}^i(T/S,\mathbb Q_\ell(j))
=F^1H^i(T,\mathbb Q_\ell(j))
  \arrow[r]
&
\operatorname{gr}^1H^i(T,\mathbb Q_\ell(j))
=E_\infty^{1,i-1}
  \arrow[d,hook]
\\
{}
&
H^1(S,R^{i-1}f_*\mathbb Q_\ell(j))
=E_2^{1,i-1}
\end{tikzcd}
\end{equation}
se denotará
\begin{equation}\tag{5.8.1.6}
U^i:\operatorname{Prim}^i(T/S,\mathbb Q_\ell(j))
\longrightarrow
H^1(S,\sisi{R^{\bullet-1}}{R^{i-1}}f_*\mathbb Q_\ell(j)).
\end{equation}
\footnote{La impresión pone un punto en la línea de base en lugar del
índice de grado de (5.8.1.6). La fuente y el blanco de (5.8.1.5)
determinan $R^{i-1}f_*$.}

\subsection*{5.8.2.}
Sea ahora $X$ proyectivo, liso e irreducible, y sea $D=\{H_\tau\}$ un
pincel de hiperplanos cuyo miembro genérico y cuyo eje cortan
transversalmente a $X$. Sea
$U\xrightarrow{\nu}\mathbb P^1$ un abierto no vacío sobre el que la
proyección $\rho:\widetilde X\to\mathbb P^1$ es lisa, y pongamos
$\widetilde X|_U=\widetilde X\times_{\mathbb P^1}U$. Se tiene
\begin{equation}\tag{5.8.2.1}
\begin{tikzcd}[column sep=large,row sep=large]
{}
&
\widetilde X
  \arrow[dl,"f"']
  \arrow[d,"\rho"]
&
\widetilde X|_U
  \arrow[l,"\alpha"']
  \arrow[d,"\rho|_U"]
\\
X
&
\mathbb P^1
&
U\arrow[l,"\nu"'].
\end{tikzcd}
\end{equation}

\pageoriginale

Como $\rho|_U:\widetilde X|_U\to U$ es propio y liso, el teorema de
especialización para la cohomología (SGA 4 XVI, 2.2) da
\begin{equation}\tag{5.8.2.2}
\operatorname{Prim}^\bullet(\widetilde X|_U/U,\mathbb Q_\ell(j))
=\ker\!\left(
H^\bullet(\widetilde X|_U,\mathbb Q_\ell(j))
\longrightarrow H^\bullet(X_{\bar\eta},\mathbb Q_\ell(j))
\right).
\end{equation}
Por consiguiente, en virtud de (5.2.3.1), se tiene:

\subsection*{Proposición 5.8.2.3.}
\underline{Bajo las hipótesis de 5.8.2, para
$x\in H^n(X,\mathbb Q_\ell(j))$, son equivalentes:}
\begin{equation}\tag{5.8.2.4}
x\in\operatorname{Prim}^n(X,\mathbb Q_\ell(j))
\quad\Longleftrightarrow\quad
\alpha^*f^*(x)\in
\operatorname{Prim}^n(\widetilde X|_U/U,\mathbb Q_\ell(j)).
\end{equation}
\underline{Si además $\rho$ satisface la hipótesis (A) de 5.3,
entonces}
\begin{equation}\tag{5.8.2.5}
x\in\operatorname{Prim}^n(X,\mathbb Q_\ell(j))
\quad\Longleftrightarrow\quad
f^*(x)\in
\operatorname{Prim}^n(\widetilde X/\mathbb P^1,\mathbb Q_\ell(j)).
\end{equation}

\noindent\underline{Demostración.}
La fórmula (5.8.2.4) ya está demostrada. Para (5.8.2.5), la hipótesis
(A) implica
\[
H^0(\mathbb P^1,R^n\rho_*\mathbb Q_\ell(j))
\xrightarrow[\nu^*]{\sim}
H^0(U,\nu^*R^n\rho_*\mathbb Q_\ell(j)).
\]
Así, \sisi{[}{}$f^*(x)$ es primitivo relativamente a
$\mathbb P^1\sisi{\alpha}{}$\sisi{]}{}
\sisi{implica que}{si y solo si}
$\alpha^*f^*(\sisi{\alpha}{x})$ es primitivo relativamente a $U$.
\footnote{La última implicación impresa contiene un $\alpha$ extraño
tras $\mathbb P^1$ y sustituye la clase $x$ por $\alpha$ en
$\alpha^*f^*(x)$. La equivalencia precedente y el isomorfismo de
restricción determinan la frase corregida.}

\subsection*{Definición 5.8.3.}
\underline{En la situación de 5.8.2, el homomorfismo compuesto}
\begin{equation}\tag{5.8.3.1}
\begin{tikzcd}[column sep=huge,row sep=huge]
\operatorname{Prim}^n(X,\mathbb Q_\ell(j))
  \arrow[r,"\alpha^*f^*"]
  \arrow[dr,"\operatorname{Grif}"']
&
\operatorname{Prim}^n(\widetilde X|_U/U,\mathbb Q_\ell(j))
  \arrow[d,"U^n"]
\\
{}
&
H^1(U,\nu^*R^{n-1}\rho_*\mathbb Q_\ell(j))
\end{tikzcd}
\end{equation}
\underline{se llama homomorfismo de Griffiths para $X$, relativo al
pincel dado $D$.}
\footnote{La tinta junto a la flecha vertical no es legible de forma
independiente. La aplicación queda, sin embargo, determinada por
(5.8.1.6): la etiqueta de grado $n$ es $U^n$.}

\pageoriginale

Añadamos al diagrama (5.8.2.1) la familia constante sobre $U$ de fibra
$X$:
\begin{equation}\tag{5.8.3.2}
\begin{tikzcd}[row sep=large,column sep=large]
{}
&
\widetilde X
  \arrow[dl,"f"']
  \arrow[dd]
&
\widetilde X|_U
  \arrow[l,"\alpha"']
  \arrow[r,"f\times\rho|_U"]
  \arrow[dd,"\rho|_U"]
&
X\times U
  \arrow[ddl,"\operatorname{pr}_2"]
\\
X&{}&{}&{}
\\
{}
&
\mathbb P^1
&
U\arrow[l,"\nu"']
&
.
\end{tikzcd}
\end{equation}

\subsection*{Teorema 5.8.4.}
\underline{Supongamos que $X$ satisface (LV) (5.2.2), y que $\rho$
satisface la hipótesis (A) (5.3). Entonces, para todo abierto no vacío
$U\xrightarrow{\nu}\mathbb P^1$ de lisura de $\rho$, el homomorfismo de
Griffiths ``módulo la parte fija'', que por definición es el compuesto}
\begin{equation}\tag{5.8.4.1}
\resizebox{\textwidth}{!}{%
\begin{tikzcd}[ampersand replacement=\&,row sep=huge,column sep=huge]
\operatorname{Prim}^n(X,\mathbb Q_\ell(j))
  \arrow[r,"\operatorname{Grif}"]
\&
H^1(U,\nu^*R^{n-1}\rho_*\mathbb Q_\ell(j))
  \arrow[d]
\\
{}
\&
H^1(U,\sisi{}{\nu^*}E^{n-1}\rho_*\sisi{(}{}\mathbb Q_\ell(j))
\simeq
\dfrac{H^1(U,\nu^*R^{n-1}\rho_*\mathbb Q_\ell(j))}
{(f\times\rho|_U)^*
 H^1(U,H^{n-1}(X,\mathbb Q_\ell(j))_U)}
\end{tikzcd}}
\end{equation}
\underline{es inyectivo.}
\footnote{La impresión omite $\nu^*$ del haz inferior en (5.8.4.1)
y deja sin cerrar su paréntesis de coeficientes. (5.8.4.3) y el mismo
nodo inferior de (5.8.4.4) dan la lectura corregida.}

\pageoriginale

\noindent\underline{Demostración.}
Por (5.5.1), sobre $\mathbb P^1$ se tiene
\begin{equation}\tag{5.8.4.2}
R^{n-1}\rho_*\mathbb Q_\ell(j)
\simeq
H^{n-1}(X,\mathbb Q_\ell(j))_{\mathbb P^1}
\oplus E^{n-1}\rho_*\mathbb Q_\ell(j),
\end{equation}
y sobre $U$,
\begin{equation}\tag{5.8.4.3}
\nu^*R^{n-1}\rho_*\mathbb Q_\ell(j)
\simeq
H^{n-1}(X,\mathbb Q_\ell(j))_U
\oplus \nu^*E^{n-1}\rho_*\mathbb Q_\ell(j).
\end{equation}
Esto permite reescribir (5.8.4.1) como
\begin{equation}\tag{5.8.4.4}
\begin{tikzcd}[row sep=huge,column sep=large]
\operatorname{Prim}^n(X,\mathbb Q_\ell(j))
  \arrow[r,"\operatorname{Grif}"]
&
H^1(U,\nu^*R^{n-1}\rho_*\mathbb Q_\ell(j))
  \arrow[d]
\\
{}
&
H^1(U,\nu^*E^{n-1}\rho_*\mathbb Q_\ell(j)).
\end{tikzcd}
\end{equation}
Por (5.7.1), la aplicación, definida gracias a la hipótesis (A),
\begin{equation}\tag{5.8.4.5}
u:\operatorname{Prim}^n(X,\mathbb Q_\ell(j))
\longrightarrow
H^1(\mathbb P^1,R^{n-1}\rho_*\mathbb Q_\ell\sisi{}{(j)})
\end{equation}
es inyectiva. Como $\mathbb P^1$ es simplemente conexa, (5.8.4.2) da
\begin{equation}\tag{5.8.4.6}
H^1(\mathbb P^1,R^{n-1}\rho_*\mathbb Q_\ell\sisi{}{(j)})
\simeq
H^1(\mathbb P^1,E^{n-1}\rho_*\mathbb Q_\ell(j)).
\end{equation}
\footnote{La impresión omite el giro de Tate $(j)$ de los blancos de
(5.8.4.5) y del miembro izquierdo de (5.8.4.6), aunque el dominio y
la descomposición (5.8.4.2) lo conservan.}
Podemos, pues, completar (5.8.4.4) en un cuadrado conmutativo.

\pageoriginale

\begin{equation}\tag{5.8.4.7}
\begin{tikzcd}[column sep=huge,row sep=huge]
\operatorname{Prim}^n(X,\mathbb Q_\ell(j))
  \arrow[r,"\operatorname{Grif}"]
  \arrow[d,hook,"u"]
&
H^1(U,\nu^*R^{n-1}\rho_*\mathbb Q_\ell(j))
  \arrow[d]
\\
H^1(\mathbb P^1,E^{n-1}\rho_*\mathbb Q_\ell(j))
  \arrow[r,"\nu^*"]
&
H^1(U,\nu^*E^{n-1}\rho_*\mathbb Q_\ell(j)).
\end{tikzcd}
\end{equation}
Solo queda ver que la flecha inferior es inyectiva. Esto resulta de
\begin{equation}\tag{5.8.4.8}
E^{n-1}\rho_*\mathbb Q_\ell(j)
\simeq
\nu_*\nu^*E^{n-1}\rho_*\mathbb Q_\ell(j)
\end{equation}
y del siguiente lema bien conocido.

\subsection*{Lema 5.8.5.}
\underline{Sea $\nu:A\to B$ un morfismo de topos, y sea $F$ un haz
sobre $A$. Entonces la aplicación canónica}
\begin{equation}\tag{5.8.5.1}
H^1(B,\nu_*F)\longrightarrow H^1(A,F)
\end{equation}
\underline{es inyectiva.}

\noindent\underline{Demostración.}
En una sucesión espectral del primer cuadrante,
\begin{equation}\tag{5.8.5.2}
E_2^{1,0}\simeq E_\infty^{1,0}.
\end{equation}
Por tanto, para la sucesión espectral de Leray
\[
E_2^{p,q}=H^p(B,R^q\nu_*F)
\Longrightarrow H^{p+q}(A,F),
\]

\pageoriginale

se tiene
\begin{equation}\tag{5.8.5.3}
E_2^{1,0}=H^1(B,\nu_*F)
\simeq E_\infty^{1,0}
=F^1H^1(A,F).
\end{equation}
Tomando $F=\nu^*E$, con $E$ un haz sobre $B$, se obtiene
\begin{equation}\tag{5.8.5.4}
H^1(B,\nu_*\nu^*E)
\lhook\joinrel\longrightarrow
H^1(A,\nu^*E).
\end{equation}

\subsection*{5.8.6.}
Sea $L/k(\eta)$ una extensión finita del cuerpo de funciones racionales
de $\mathbb P^1$, y pongamos
\begin{equation}\tag{5.8.6.1}
S=\text{la normalización de $\mathbb P^1$ en $L$}.
\end{equation}
Hagamos el cambio de base $S\to\mathbb P^1$ en (5.8.3.2):
\begin{equation}\tag{5.8.6.1 bis}
\begin{tikzcd}[column sep=huge,row sep=large]
{} & \widetilde X \arrow[dl,"f"] \arrow[dd,"\rho"]
& \widetilde X_S \arrow[l,"\alpha_S"'] \arrow[dd,"\rho_S"] \\
X&{}&{} \\
{} & \mathbb P^1 & S\arrow[l,"\pi"'] .
\end{tikzcd}
\end{equation}
Para toda inmersión abierta $V\xhookrightarrow{\mu}S$ sobre la que
$\rho_S$ es lisa, el cambio de base por
$V\xrightarrow{\pi\mu}\mathbb P^1$ en (5.8.3.2) da
\begin{equation}\tag{5.8.6.2}
\begin{tikzcd}[column sep=huge,row sep=large]
{} & \widetilde X \arrow[dl,"f"] \arrow[dd,"\rho"]
& \widetilde X_V \arrow[l,"\alpha_V"'] \arrow[dd,"\rho_V"]
  \arrow[dr,"f\times\rho_V"] &{} \\
X&{}&{}&X\times V\arrow[dl,"\operatorname{pr}_2"] \\
{} & \mathbb P^1 & V\arrow[l,"\pi\mu"'] &{} .
\end{tikzcd}
\end{equation}

\pageoriginale

Se sigue teniendo (cf. (5.8.2.3))
\begin{equation}\tag{5.8.6.3}
\alpha_V^*f^*\operatorname{Prim}^n(X,\mathbb Q_\ell(j))
\subset
\operatorname{Prim}^n(\widetilde X_V/V,\mathbb Q_\ell(j)),
\end{equation}
de modo que puede hablarse de nuevo del homomorfismo de Griffiths
\begin{equation}\tag{5.8.6.4}
\operatorname{Prim}^n(X,\mathbb Q_\ell(j))
\xrightarrow{\operatorname{Grif}}
H^1(V,R^{n-1}\rho_{V*}\mathbb Q_\ell(j)).
\end{equation}

\subsection*{Corolario 5.8.7.}
\underline{Con la notación precedente y bajo las hipótesis de 5.8.4,
el homomorfismo de Griffiths ``módulo la parte fija'', es decir, el
compuesto}
\begin{equation}\tag{5.8.7.1}
\resizebox{\textwidth}{!}{%
\begin{tikzcd}[ampersand replacement=\&,column sep=huge,row sep=huge]
\operatorname{Prim}^n(X,\mathbb Q_\ell(j))
  \arrow[r,"\operatorname{Grif}"]
\& H^1(V,R^{n-1}\rho_{V*}\mathbb Q_\ell(j)) \arrow[d] \\
{}
\& H^1(V,(\pi\mu)^*E^{n-1}\rho_*\mathbb Q_\ell(j))
\simeq
\dfrac{H^1(V,R^{n-1}\rho_{V*}\mathbb Q_\ell(j))}
{(f\times\rho_V)^*H^1(V,H^{n-1}(X,\mathbb Q_\ell(j))_V)}
\end{tikzcd}}
\end{equation}
\underline{es inyectivo.}

\pageoriginale

\noindent\underline{Demostración.}
Por el teorema de especialización para la cohomología,
\[
R^{n-1}\rho_{V*}\mathbb Q_\ell(j)
=(\pi\mu)^*R^{n-1}\rho_*\mathbb Q_\ell(j).
\]
Así, la descomposición (5.8.4.2) reescribe (5.8.7.1) como
\[
\begin{tikzcd}[column sep=huge,row sep=huge]
\operatorname{Prim}^n(X,\mathbb Q_\ell(j))
  \arrow[r,"\operatorname{Grif}"]
& H^1(V,(\pi\mu)^*R^{n-1}\rho_*\mathbb Q_\ell(j)) \arrow[d] \\
{}
& H^1(V,(\pi\mu)^*E^{n-1}\rho_*\mathbb Q_\ell(j)).
\end{tikzcd}
\]
Basta, pues, demostrar la inyectividad de
\begin{equation}\tag{5.8.7.3}
\operatorname{Prim}^n(X,\mathbb Q_\ell(j))
\longrightarrow
H^1(V,\mu^*\pi^*E^{n-1}\rho_*\mathbb Q_\ell(j))
\end{equation}
para inmersiones abiertas $V\xhookrightarrow{\mu}S$ arbitrariamente
pequeñas. Nos reducimos al caso
\[
\left\{
\begin{aligned}
U&\xhookrightarrow{\nu}\mathbb P^1
   &&\text{abierto sobre el que $\rho$ es lisa},\\
V&\xrightarrow{\pi_U}U
   &&\text{normalización de $U$ en $L$}.
\end{aligned}
\right.
\]
Del cuadrado cartesiano
\[
\begin{tikzcd}[column sep=huge,row sep=huge]
V\arrow[r,hook,"\mu"]\arrow[d,"\pi_U"']
&S\arrow[d,"\pi"]\\
U\arrow[r,hook,"\nu"]
&\mathbb P^1
\end{tikzcd}
\]

\pageoriginale

se obtiene
\[
\mu^*\pi^*E^{n-1}\rho_*\mathbb Q_\ell(j)
=\pi_U^*\nu^*E^{n-1}\rho_*\mathbb Q_\ell\sisi{}{(j)}.
\]
Por tanto, (5.8.7.3) se factoriza como
\[
\begin{tikzcd}[column sep=huge,row sep=huge]
\operatorname{Prim}^n(X,\mathbb Q_\ell(j)) \arrow[r]
&H^1(U,\nu^*E^{n-1}\rho_*\mathbb Q_\ell(j))
  \arrow[d,"\pi_U^*"] \\
{}
&H^1(V,\pi_U^*\nu^*E^{n-1}\rho_*\mathbb Q_\ell(j)).
\end{tikzcd}
\]
Por (5.8.4), basta verificar que
\[
\pi_U^*:H^1(U,\nu^*E^{n-1}\rho_*\mathbb Q_\ell(j))
\longrightarrow
H^1(V,\pi_U^*\nu^*E^{n-1}\rho_*\mathbb Q_\ell(j))
\]
es inyectiva. Como $\pi_U:V\to U$ es finita y plana, la teoría de la
traza [SGA 4 XVII, 6.2] proporciona, para todo $\mathbb Q_\ell$-haz
$F$ sobre $U$, por ejemplo
$F=\nu^*E^{n-1}\rho_*\mathbb Q_\ell(j)$, una aplicación
\[
\operatorname{tr}_{\pi_U}:H^1(V,\pi_U^*\sisi{}{F})
\longrightarrow H^1(U,F\sisi{(j)}{})
\]
tal que
\[
\sisi{\operatorname{tr}_{\pi}\circ\pi^*}
{\operatorname{tr}_{\pi_U}\circ\pi_U^*}
=\deg(\pi_U)\,\operatorname{id}.
\]
Se sigue que $H^*(U,F)\to H^*(V,\pi_U^*F)$ es inyectiva, y termina la
demostración.
\footnote{La impresión omite el giro de Tate $(j)$ en una fórmula de cambio
de base y la transcripción congelada corrompe la fuente y el blanco de
la traza. Las fórmulas mostradas son las identidades estándar de
imagen inversa y traza finita plana utilizadas en la conclusión.}

\pageoriginale

\subsection*{5.9. Un ejemplo en el que no se satisface la hipótesis (A)}

Sea $X$ una superficie cuádrica lisa de $\mathbb P^3$. Entonces $X$
satisface (LV) (5.2.4). Tomemos un pincel, de Lefschetz si se quiere,
de \underline{hiperplanos} cuyo miembro genérico y cuyo eje cortan
transversalmente a $X$. Entonces, como me hizo observar Deligne,
\underline{la proyección}
\[
\rho:\widetilde X\longrightarrow\mathbb P^1
\]
\underline{no satisface (A)}. En efecto, $R^2\rho_*\mathbb Q_\ell$ no
es constante: las fibras singulares son cónicas singulares y, por
tanto, reducibles; en un pincel de Lefschetz cada una es la unión de
dos rectas que se cortan en un punto, que es la definición en este
caso sencillo. Así,
\[
\dim_{\mathbb Q_\ell}H^2(\text{fibra singular},\mathbb Q_\ell)=2,
\]
mientras que
\[
\dim_{\mathbb Q_\ell}H^2(\text{fibra lisa},\mathbb Q_\ell)=1.
\]

Además, aquí es falsa la conclusión del teorema 5.8.4. Sea
$U\xhookrightarrow{\nu}\mathbb P^1$ un abierto no vacío sobre el que
$\rho$ es lisa. Como las fibras de $\rho$ sobre $U$ son curvas lisas de
género cero,
\[
\nu^*R^1\rho_*\mathbb Q_\ell=0
\]
y, por tanto,
\[
H^1(U,\nu^*R^1\rho_*\mathbb Q_\ell)=0.
\]
En consecuencia, el homomorfismo de Griffiths
\[
\operatorname{Prim}^2(X,\mathbb Q_\ell)
\xrightarrow{\operatorname{Grif}}
H^1(U,\nu^*R^1\rho_*\mathbb Q_\ell)
\]
es nulo. Pero, sobre una superficie cuádrica,
\[
\dim_{\mathbb Q_\ell}\operatorname{Prim}^2(X,\mathbb Q_\ell)=1.
\]
Asimismo falla aquí la conclusión de 5.7.

\pageoriginale

\section*[6. Ciclos evanescentes (Picard--Lefschetz)]{\texorpdfstring{6. Aplicación de la teoría local de los ciclos evanescentes
(Picard--Lefschetz)}{6. Aplicación de la teoría local de los ciclos evanescentes (Picard--Lefschetz)}}

\subsection*{Proposición 6.1.}

Sea $Z$ una hipersuperficie del espacio proyectivo $\mathbb P^r_k$,
$r\geq2$, donde $k$ es un cuerpo algebraicamente cerrado. Sea
\[
i_o:D_o\lhook\joinrel\longrightarrow\mathbb P^r
\]
la inclusión de una recta que no corta a $Z$ más que en su lugar liso
y que la corta transversalmente. Elijamos un punto geométrico $\bar a$
de $D_o-D_o\cap Z$, y denotemos por
\[
\pi_1^{\mathrm{mod}}(D_o-D_o\cap Z,\bar a)
\]
el cociente de $\pi_1(D_o-D_o\cap Z,\bar a)$ que clasifica los
recubrimientos étale de $D_o-D_o\cap Z$ moderadamente ramificados en
todo punto de $D_o\cap Z$. Entonces:

\pageoriginale

\subsubsection*{6.1.1.}

El homomorfismo canónico
\begin{equation}\tag{6.1.1.1}
i_{o*}:\pi_1^{\mathrm{mod}}(D_o-D_o\cap Z,\bar a)
\longrightarrow
\pi_1^{\mathrm{mod}}(\mathbb P^r-Z,\bar a)
\end{equation}
es sobreyectivo.

\subsubsection*{6.1.2.}

Toda elección de un punto geométrico
$\bar x_o\in D_o\cap Z$ y de un ``camino'' en $D_o$ desde $\bar x_o$
hasta $\bar a$ determina un homomorfismo
\begin{equation}\tag{6.1.2.1}
\prod_{\ell\ne\operatorname{char}(k)}\mathbb Z_\ell(1)
\longrightarrow
\pi_1^{\mathrm{mod}}(D_o-D_o\cap Z,\bar a),
\end{equation}
cuya imagen es ``el'' grupo de inercia en $\bar x_o$. Al componer este
morfismo con $i_{o*}$ se obtiene un homomorfismo
\begin{equation}\tag{6.1.2.2}
\prod_{\ell\ne\operatorname{char}(k)}\mathbb Z_\ell(1)
\longrightarrow
\pi_1^{\mathrm{mod}}(\mathbb P^r-Z,\bar a).
\end{equation}
Si $Z$ es irreducible, todos estos homomorfismos son conjugados entre
sí.

\subsubsection*{Demostración.}

Sea $T$ el subesquema de la grassmanniana formado por las rectas de
$\mathbb P^r$ que
\begin{enumerate}[label=(\roman*)]
\item pasan por $\bar a$;
\item no cortan a $Z$ más que en su lugar liso, y la cortan allí
  transversalmente.
\end{enumerate}
Sea $\eta$ su punto genérico y sea
\[
i_\eta:D_\eta\longrightarrow\mathbb P^r\otimes_k k(\eta)
\]
la recta correspondiente (por tanto, una recta que ``une $\bar a$ con
el punto genérico de $\mathbb P^r_k$''). Elijamos asimismo un punto
geométrico $\bar\eta$ sobre $\eta$, y denotemos por $D_{\bar\eta}$ la
recta ``geométrica'' correspondiente. Sea $o\in T$ el punto
correspondiente a la recta dada $D_o$.

\pageoriginale

Por [SGA 1, XIII 2.8], existe un homomorfismo de especialización
(dependiente de un camino en $T$ desde $\bar\eta$ hasta $o$, y por
tanto definido salvo conjugación)
\begin{equation}\tag{6.1.3}
\pi_1^{\mathrm{mod}}(D_{\bar\eta}-D_{\bar\eta}\cap Z,\bar a)
\xrightarrow{\ \bar\eta\longrightarrow o\ }
\pi_1^{\mathrm{mod}}(D_o-D_o\cap Z,\bar a)
\end{equation}
con las propiedades siguientes.

\subsubsection*{6.1.4.}

Los homomorfismos de monodromía local (6.1.2.1)
\[
\prod_{\ell\ne\operatorname{char}(k)}\mathbb Z_\ell(1)
\longrightarrow
\pi_1^{\mathrm{mod}}(D_o-D_o\cap Z,\bar a)
\]
se obtienen componiendo la flecha de especialización (6.1.3) con los
homomorfismos de monodromía local para $D_{\bar\eta}$:
\begin{equation}\tag{6.1.4.1}
\prod_{\ell\ne\operatorname{char}(k)}\mathbb Z_\ell(1)
\longrightarrow
\pi_1^{\mathrm{mod}}(D_{\bar\eta}-D_{\bar\eta}\cap Z,\bar a)
\xrightarrow{\ \bar\eta\longrightarrow o\ }
\pi_1^{\mathrm{mod}}(D_o-D_o\cap Z,\bar a).
\end{equation}

\subsubsection*{6.1.5.}

El diagrama siguiente es conmutativo:
\begingroup
\small
\begin{equation}\tag{6.1.5}
\begin{tikzcd}[column sep=large,row sep=large]
{} & \pi_1^{\mathrm{mod}}(\mathbb P^r-Z,\bar a) & {} \\
\pi_1^{\mathrm{mod}}(D_\eta-D_\eta\cap Z,\bar a)
  \arrow[ur,"i_{\eta*}"]
& {}
& \pi_1^{\mathrm{mod}}(D_o-D_o\cap Z,\bar a)
  \arrow[ul,"i_{o*}"'] \\
{} & \pi_1^{\mathrm{mod}}(D_{\bar\eta}-D_{\bar\eta}\cap Z,\bar a)
  \arrow[uu,"i_{\bar\eta*}"']
  \arrow[ul]
  \arrow[ur,"\bar\eta\longrightarrow o"'] & {} .
\end{tikzcd}
\end{equation}
\endgroup

\pageoriginale

Admitamos esto por el momento y demostremos primero 6.1.1. Por la
conmutatividad de 6.1.5 basta probar que
\[
i_{\bar\eta*}:
\pi_1^{\mathrm{mod}}(D_{\bar\eta}-D_{\bar\eta}\cap Z,\bar a)
\longrightarrow
\pi_1^{\mathrm{mod}}(\mathbb P^r-Z,\bar a)
\]
es sobreyectivo. Para ello basta demostrar que el homomorfismo análogo
de los grupos fundamentales ordinarios
\begin{equation}\tag{6.1.6}
\sisi{i_{\bar\eta}}{i_{\bar\eta*}}:
\pi_1(D_{\bar\eta}-D_{\bar\eta}\cap Z,\bar a)
\longrightarrow
\pi_1(\mathbb P^r-Z,\bar a)
\end{equation}
es sobreyectivo. En otras palabras, para todo recubrimiento étale
$E\to\mathbb P^r-Z$ con $E$ conexo (y por tanto irreducible, pues
$\mathbb P^r-Z$ es liso y, en consecuencia, unirrama
[EGA IV, 18.10.3]), su imagen inversa
\[
E_{\bar\eta}
=E\times_{\mathbb P^r-Z}(D_{\bar\eta}-D_{\bar\eta}\cap Z)
\]
sobre $D_{\bar\eta}-D_{\bar\eta}\cap Z$ debe seguir siendo
geométricamente irreducible. Basta aplicar la variante siguiente del
teorema de Bertini ([EGA V, referencia dejada en blanco en la fuente])
al compuesto
$E\to\mathbb P^r-Z\lhook\joinrel\longrightarrow\mathbb P^r$ y a $s=1$.

\subsubsection*{6.1.6.1. Teorema de Bertini.}

Sea $k$ un cuerpo, sea $\mathbb P^r=\mathbb P^r_k$ y sea
$f:X\to\mathbb P^r$ un morfismo no constante, donde $X$ es un esquema
algebraico geométricamente irreducible. Para cada entero
$s<\dim\overline{f(X)}$, denotemos por $L_t$ el subespacio lineal
\underline{genérico} de codimensión $s$ de $\mathbb P^r$. Entonces
$X_t=f^{-1}(L_t)$ es geométricamente irreducible.

Demostremos ahora 6.1.2. Por 6.1.4 y la conmutatividad de (6.1.5), basta
probar que todos los homomorfismos de monodromía local (6.1.3)
\underline{para} $D_{\bar\eta}$,

\pageoriginale

\begin{equation}\tag{6.1.7.1}
\prod_{\ell\ne\operatorname{char}(k)}\mathbb Z_\ell(1)
\longrightarrow
\pi_1^{\mathrm{mod}}(D_{\bar\eta}-D_{\bar\eta}\cap Z,\bar a),
\end{equation}
se vuelven conjugados tras componerlos con el homomorfismo
\begin{equation}\tag{6.1.7.2}
\pi_1^{\mathrm{mod}}(D_{\bar\eta}-D_{\bar\eta}\cap Z,\bar a)
\longrightarrow
\pi_1^{\mathrm{mod}}(D_\eta-D_\eta\cap Z,\bar a).
\end{equation}
Para probarlo, observemos que, para toda extensión de Galois finita
$K$ de $k(\eta)$, con
\[
k(\eta)\subset K\subset k(\bar\eta),
\]
el esquema
\[
D_K\stackrel{\mathrm{def}}{=}D_\eta\times_{k(\eta)}K
\]
es un recubrimiento étale de $D_\eta$. Por tanto,
$D_K-D_K\cap Z$ es un recubrimiento moderado de
$D_\eta-D_\eta\cap Z$, y
$\pi_1^{\mathrm{mod}}(D_\eta-D_\eta\cap Z,\bar a)$ actúa sobre el
esquema $D_K\cap Z$. Tomamos $K$ suficientemente grande para que todo
punto de $D_K\cap Z$ sea $K$-racional; esto es posible porque, por
hipótesis, $D_\eta\cap Z$ es étale sobre $k(\eta)$.

Queda demostrar que
$\pi_1^{\mathrm{mod}}(D_\eta-D_\eta\cap Z,\bar a)$ actúa
\underline{transitivamente} sobre los puntos de $D_K\cap Z$. Esto
probará que los homomorfismos compuestos de inercia local
\[
\prod_{\ell\ne\operatorname{char}(k)}\mathbb Z_\ell(1)
\longrightarrow
\pi_1^{\mathrm{mod}}(D_\eta-D_\eta\cap Z,\bar a)
\]
asociados a los distintos puntos de $D_K\cap Z$ son conjugados. Ahora
bien, la acción sobre $D_K\cap Z$ se factoriza por el cociente canónico
\[
\pi_1^{\mathrm{mod}}(D_\eta-D_\eta\cap Z,\bar a)
\longrightarrow
\pi_1(\eta,\bar\eta)
=\sisi{\operatorname{Gal}(\,/k(\eta))}
{\operatorname{Gal}\bigl(k(\bar\eta)/k(\eta)\bigr)}.
\]

\pageoriginale

La acción de
$\operatorname{Gal}\bigl(k(\bar\eta)/k(\eta)\bigr)$ sobre
$D_K\cap Z$ es la compatible con la inclusión
\[
D_K\cap Z\lhook\joinrel\longrightarrow\mathbb P^r_k\times_k K,
\]
donde el grupo de Galois actúa sobre $\mathbb P^r_k\times_k K$
únicamente por el segundo factor. Se sigue que la afirmación deseada
equivale a que $D_\eta\cap Z$ tenga un solo punto, es decir, a que sea
el espectro de un \underline{cuerpo}. En términos elementales, esto no
es sino el hecho siguiente, de tipo Bertini (se toma para $Z$ una
ecuación afín $F=0$ y se supone que $\bar a$ es el origen).

\subsubsection*{Hecho 6.1.8.}

Sea $k$ un cuerpo algebraicamente cerrado, sea $r\ge2$ y sea
$F(X_1,\ldots,X_r)\in k[X_1,\ldots,X_r]$ un polinomio irreducible.
Entonces
\[
G(T)=F(X_1T,\ldots,X_rT)\in k(X_1,\ldots,X_r)[T]
\]
es irreducible como polinomio en $T$ con coeficientes en
$k(X_1,\ldots,X_r)$.

Esto termina la demostración de 6.1.2.

\subsubsection*{Observación 6.1.9.}

Otro método para demostrar 6.1.2, válido bajo condiciones
considerablemente más generales, consistiría simplemente en establecer
una compatibilidad general para los ``homomorfismos de monodromía'':
los homomorfismos compuestos de 6.1.2 son conjugados del homomorfismo de
monodromía relativo al punto genérico de $Z$, y por tanto son conjugados
entre sí.

\pageoriginale

\subsubsection*{Corolario 6.1.10.}

Sea $X$ una variedad proyectiva, lisa y conexa, de dimensión $n$ sobre
un cuerpo algebraicamente cerrado $k$. Supongamos que
$n=\dim X$ es par, o que
$\sisi{\operatorname{card}(k)}{\operatorname{char}(k)}\ne2$. Sea $D$
un pincel de Lefschetz para $X$, sea $\eta$ su punto genérico y sea
$\bar\eta$ un punto geométrico sobre $\eta$. Entonces la representación
natural de $\pi_1(D-D\cap\check X,\bar\eta)$ sobre
$H^{n-1}(X_{\bar\eta},\mathbb Q_\ell)$ se factoriza por
$\pi_1^{\mathrm{mod}}(D-D\cap\check X,\bar\eta)$. Si esta
representación no es trivial, entonces $D\cap\check X\ne\varnothing$ y
el grupo de inercia en todo punto de $D\cap\check X$ actúa no
trivialmente.

\subsubsection*{Demostración.}

Si $\dim(\check X)\le r-2$, entonces $D\cap\check X$ es vacío y no hay
nada que demostrar, pues $\mathbb P^1$ es simplemente conexa. En caso
contrario, $\check X$ es una hipersuperficie de $\mathbb P^r$
[XVII, 3.1.4], y 6.1.1 se aplica por 3.2.10 y XVII 3.5, tomando
$Z=\check X$ y $D_o=D$. El hecho de que la acción de $\pi_1$ sobre
$H^{n-1}(X_{\bar\eta},\mathbb Q_\ell)$ sea moderada se estableció en la
teoría local mediante la fórmula de Picard--Lefschetz [XV, 3.4], que
recordaremos más abajo. Finalmente, el hecho no trivial de que
$\pi_1^{\mathrm{mod}}(D-D\cap\check X,\bar\eta)$ está engendrado por
los grupos de inercia en los distintos puntos de $D\cap\check X$
[SGA 1, XIII 2.10] da la conclusión.

Utilizaremos ahora el resultado precedente para estudiar la condición
(A) de (5.3.5). Recordemos primero la consecuencia siguiente de la
fórmula de Picard--Lefschetz [XV, 3.4].

\pageoriginale

\subsubsection*{Proposición 6.2.}

Sea $X$ una variedad proyectiva, lisa y conexa, de dimensión $n$ sobre
el cuerpo algebraicamente cerrado $k$. Sea $D$ un pincel de Lefschetz,
sea $\eta$ su punto genérico y sea $\bar\eta$ un punto geométrico sobre
$\eta$. Entonces el pincel $D$ satisface la condición (A) si y solo si
el grupo de inercia en todo punto de $D\cap\check X$ actúa no
trivialmente sobre $H^{n-1}(X_{\bar\eta},\mathbb Q_\ell)$.

\subsubsection*{Teorema 6.3.}

Bajo las hipótesis de 6.2 se tiene:

\medskip
\noindent
\textup{(6.3.1)} Si $n=\dim X$ es impar, todo pincel de Lefschetz para
$X$ satisface la condición (A).

\medskip
\noindent
\textup{(6.3.2)} Supongamos que $n=\dim X$ es par, o que
$\operatorname{char}(k)\ne2$. Entonces son equivalentes las condiciones
siguientes:
\begin{enumerate}
\item[(i)] $D$ no satisface la condición (A) de (5.3.5);
\item[(ii)] $\pi_1(D-D\cap\check X)$ actúa trivialmente sobre
  $H^{n-1}(X_{\bar\eta},\mathbb Q_\ell)$, y
  $D\cap\check X\ne\varnothing$;
\item[(iii)] existe un punto de $D\cap\check X$ cuyo grupo de inercia
  actúa trivialmente sobre $H^{n-1}(X_{\bar\eta},\mathbb Q_\ell)$;
\item[(iv)] el haz $R^{n-1}\rho_*\mathbb Q_\ell$ es constante sobre
  $D$, y $D\cap\check X$ no es vacío;
\item[(v)] la restricción
  $R^{n-1}\rho_*\mathbb Q_\ell\vert_U$ es constante sobre $U$, donde
  $U$ es un abierto no vacío de $D$ sobre el que $\rho$ es lisa, y
  $\rho$ no es lisa.
\end{enumerate}

\noindent
\textup{(6.3.3)} Si, además, $X$ satisface la condición (LV) de
(5.2.2), las condiciones (i)--(v) son asimismo equivalentes a:
\begin{enumerate}
\item[(vi)] la aplicación de restricción
\[
H^{n-1}(X,\mathbb Q_\ell)
\longrightarrow H^{n-1}(X_{\bar\eta},\mathbb Q_\ell)
\]
es sobreyectiva (y por tanto un isomorfismo por el teorema de Lefschetz
débil).
\end{enumerate}

\subsubsection*{Demostración.}

La afirmación 6.3.1 resulta de XV 3.4. Para las equivalencias de 6.3.2,
(i) equivale a (iii) por 6.2, (ii) a (iii) por 6.1.6, (ii) a (v) por
tautología y (iv) a (v) por 5.3.4. Finalmente, si $X$ satisface la
condición (LV), 5.6.10 da la equivalencia de (ii) y (vi).

\pageoriginale

\subsubsection*{Observación 6.3.4.}

La condición (A) no depende de la elección del pincel de Lefschetz
$D$, sino únicamente de la inmersión (de Lefschetz)
$i:X\longrightarrow\mathbb P^r$ considerada. Esto es claro por (6.3.1)
si $\dim X$ es impar. Si $\dim X$ es par, la condición (A) significa
también que el sistema local sobre $\mathbb P^r-\check X$ formado por
los $H^{n-1}(X\cap H_s,\mathbb Q_\ell)$ es constante. Esto resulta del
criterio (ii), del hecho de que el sistema local considerado está
moderadamente ramificado en el punto genérico de $\check X$ y de la
sobreyectividad de (6.1.2).

\subsubsection*{Corolario 6.4.}

Bajo las hipótesis de 6.2, supongamos
\sisi{que $n=\dim X$ es par, o que $\operatorname{char}(k)\ne2$ y que
$X$ satisface la condición (LV)}
{que ($n=\dim X$ es par o $\operatorname{char}(k)\ne2$) y que $X$
satisface la condición (LV)}.
Entonces existe un entero $d_o$ tal que, para $d\ge d_o$, todo pincel de

\pageoriginale

Lefschetz de secciones de $X$ por hipersuperficies de grado $d$
satisface la condición (A).

\subsubsection*{Demostración.}

Por la equivalencia (i)$\Longleftrightarrow$(vi) de (6.3.2), basta
probar que, si $Y(d)$ denota una sección hipersuperficie lisa de $X$ de
grado $d$, entonces
\begin{equation}\tag{6.4.1}
\dim H^{n-1}(Y(d),\mathbb Q_\ell)\longrightarrow\infty
\qquad\text{cuando }d\longrightarrow\infty.
\end{equation}
En efecto, se tiene el resultado preciso siguiente.

\subsubsection*{Lema 6.4.2.}

Sea $X$ una variedad proyectiva, lisa y conexa, de dimensión $n\ge1$
sobre un cuerpo algebraicamente cerrado $k$, y denotemos por $Y(d)$ una
sección hipersuperficie de grado $d$. Entonces existe un polinomio
$f(T)\in\mathbb Z[T]$ tal que
\begin{align}
f(T)&=\deg(X)T^n+\text{un polinomio de grado a lo sumo }n-1,
  \tag{6.4.2.1}\\
f(d)&=\dim H^{n-1}(Y(d),\mathbb Q_\ell)\qquad(d\ge1).
  \tag{6.4.2.2}
\end{align}

\subsubsection*{Demostración.}

Por XVII 5.7.5, si $C(X)$ denota la clase de Chern total de $X$ y $H$
la primera clase de Chern de un hiperplano, la característica de
Euler--Poincaré de $Y(d)$ es
\begin{equation}\tag{6.4.2.3}
\sisi{\chi(Y(d)}{\chi(Y(d))}
=\int_X\frac{dH\,C(X)}{1+dH}
=\sum_{i=1}^{n}(-1)^{i-1}d^i\int_X H^iC_{n-i}.
\end{equation}

\pageoriginale

Por otra parte, la dualidad de Poincaré y Lefschetz débil dan
\begin{equation}\tag{6.4.2.4}
\chi(Y(d))
=(-1)^{n-1}\dim H^{n-1}(Y(d),\mathbb Q_\ell)
+2\sum_{i\ge1}(-1)^{n-1-i}
  \dim H^{n-1-i}(X,\mathbb Q_\ell).
\end{equation}
Al combinar (6.4.2.3) y (6.4.2.4) se obtiene el resultado.

\subsection*{6.5. Recordatorio de la fórmula de Picard--Lefschetz
(XV 3.4)}

Sea $S$ el espectro de un anillo de valoración estrictamente
henseliano, con punto genérico $\eta$, punto genérico geométrico
$\bar\eta$ y punto cerrado $s$. Sea
\[
\rho:X\longrightarrow S
\]
un morfismo propio, donde $X$ es un esquema regular conexo de dimensión
$n$. Supongamos que $\rho$ es liso salvo en un único punto $x$ de la
fibra especial, y que $x$ es una singularidad cuadrática ordinaria de
dicha fibra.

Denotemos por $I$ el grupo de inercia de $S$ en $s$. Por XV 3.4, $I$
actúa trivialmente sobre $H^i(X_{\bar\eta},\mathbb Q_\ell)$ para
$i\ne n-1$, y su acción sobre
$H^{n-1}(X_{\bar\eta},\mathbb Q_\ell)$ se describe como sigue.

\subsubsection*{6.5.1.}

Si $n=2m+1$, entonces $I$ actúa a través de un cociente isomorfo a
$\mathbb Z/2\mathbb Z$ (si $\operatorname{char}(k)\ne2$, existe un solo
cociente de este tipo). Existe un elemento
\[
\delta\in H^{n-1}(X_{\bar\eta},\mathbb Q_\ell(m)),
\]
determinado de manera única salvo signo y llamado el ciclo
\underline{evanescente en $s$}, tal que

\pageoriginale

\begin{equation}\tag{6.5.1.1}
\delta\cdot\delta=(-1)^m2
\quad\text{en}\quad
H^{2n-2}(X_{\bar\eta},\mathbb Q_\ell(n-1))\simeq\mathbb Q_\ell.
\end{equation}
El elemento no trivial $\sigma$ de $\mathbb Z/2\mathbb Z$ actúa sobre
$H^{n-1}(X_{\bar\eta},\mathbb Q_\ell)$ mediante la simetría
\begin{equation}\tag{\sisi{6.5.2.2}{6.5.1.2}}
\sigma x=x-(-1)^m(x\cdot\delta)\delta.
\end{equation}
Aquí $x\cdot\delta$ pertenece a
$H^{2n-2}(X_{\bar\eta},\mathbb Q_\ell(m))\simeq\mathbb Q_\ell(-m)$,
de modo que $(x\cdot\delta)\delta$ define un elemento de
$H^{n-1}(X_{\bar\eta},\mathbb Q_\ell)$.

\subsubsection*{6.5.2.}

Si $n=2m+2$, entonces $I$ actúa a través de su cociente canónico
$\mathbb Z_\ell(1)$. Existe un elemento
\[
\delta\in
\sisi{H^{n-1}(X_{\bar\eta},\mathbb Q_\ell(m)}
{H^{n-1}(X_{\bar\eta},\mathbb Q_\ell(m))},
\]
determinado canónicamente salvo signo y llamado el ciclo
\underline{evanescente en $s$}, tal que
$\sigma\in\mathbb Z_\ell(1)$ actúa sobre
$x\in H^{n-1}(X_{\bar\eta},\mathbb Q_\ell)$ mediante la transvección
\begin{equation}\tag{6.5.2.1}
\sigma(x)=x-(-1)^m\sigma\,(x\cdot\delta)\delta.
\end{equation}
En efecto, $x\cdot\delta$ pertenece a
$H^{2n-2}(X_{\bar\eta},\mathbb Q_\ell(m))
\simeq\mathbb Q_\ell(-m-1)$. Su multiplicación por
$\sigma\in\mathbb Z_\ell(1)$ da, pues, un elemento
$\sigma(x\cdot\delta)\in\mathbb Q_\ell(-m)$, cuyo producto por
$\delta\in H^{n-1}(X_{\bar\eta},\mathbb Q_\ell(m))$ pertenece a
$H^{n-1}(X_{\bar\eta},\mathbb Q_\ell)$.

Volvamos a la situación global.

\subsubsection*{Teorema 6.6.}

Sea $X$ una variedad proyectiva, lisa y conexa, de dimensión $n$ sobre
un cuerpo algebraicamente cerrado $k$. Supongamos que $n$ es par o que
$\operatorname{char}(k)\ne2$. Sea $D$ un pincel de Lefschetz para $X$,
sea $\eta$ su punto genérico, sea $\bar\eta$ un punto geométrico sobre
$\eta$, y sean $s_1,\ldots,s_N$ los puntos de $D\cap\check X$.

\pageoriginale

Para cada $i$, elijamos un isomorfismo $\ell_i$ entre $\bar\eta$ y el
punto genérico geométrico de $S_i$, el henselizado estricto de $D$ en
$s_i$. Por medio de $\ell_i$, sea $I_i$ el grupo de inercia
\sisi{de $\delta_i$ en $S_i$}{en $s_i$}, que actúa sobre
$H^{n-1}(X_{\bar\eta},\mathbb Q_\ell(j))$, y sea $\delta_i$ el ciclo
evanescente en $s_i$ definido en 6.5. Así,
\[
\delta_i\in
\sisi{H^{n-1}(X,\mathbb Q_\ell(m)}
{H^{n-1}(X_{\bar\eta},\mathbb Q_\ell(m))},
\qquad
m=\begin{cases}
(n-1)/2,& n\text{ impar},\\
n/2-1,& n\text{ par}.
\end{cases}
\]
Elijamos los $\ell_i$ de modo que los subgrupos
\[
I_i\subset
\pi=\pi_1^{\mathrm{mod}}(D-D\cap\check X,\bar\eta)
\]
engendren topológicamente $\pi$, como permite SGA 1, XIII 2.10.
Entonces:

\medskip
\noindent
\textup{6.6.1.} Si $X$ satisface la condición (LV), el espacio de
cohomología evanescente
\[
\sisi{E^{n-1}(X_{\bar\eta},\mathbb Q_\ell(m)}
{E^{n-1}(X_{\bar\eta},\mathbb Q_\ell(m))}
\stackrel{\mathrm{def}}{=}
H^{n-1}(X,\mathbb Q_\ell(m))^\perp
\subset H^{n-1}(X_{\bar\eta},\mathbb Q_\ell(m))
\]
(cf. 5.4.4) está engendrado por los ciclos evanescentes
$\delta_i$, $i=1,\ldots,N$.

\medskip
\noindent
\textup{6.6.2.} Los ciclos evanescentes $\delta_i$ son conjugados,
salvo signo, bajo el grupo de monodromía global
$\pi_1(D-D\cap\check X)$.

\subsubsection*{Demostración.}

Para 6.6.1, la fórmula de Picard--Lefschetz 6.5 y el hecho de que los
$I_i$ engendran $\pi=\pi_1(D-D\cap\check X)$ dan

\pageoriginale

\begin{equation}\tag{6.6.3}
H^{n-1}(X_{\bar\eta},\mathbb Q_\ell(m))^\pi
=\bigl\langle\delta_1,\ldots,\delta_N\bigr\rangle^\perp
\quad\text{en }H^{n-1}(X_{\bar\eta},\mathbb Q_\ell(m)).
\end{equation}
Por 5.6.10,
\begin{equation}\tag{6.6.4}
H^{n-1}(X_{\bar\eta},\mathbb Q_\ell(m))^\pi
=\operatorname{Im}\bigl(
H^{n-1}(X,\mathbb Q_\ell(m))
\longrightarrow H^{n-1}(X_{\bar\eta},\mathbb Q_\ell(m))
\bigr).
\end{equation}
Al tomar ortogonales en (6.6.3) se obtiene 6.6.1.

Para 6.6.2, observemos simplemente que el ciclo evanescente $\delta_i$
está determinado salvo signo por la representación del grupo de
inercia en $s_i$. Apliquemos 6.1.2, que da la conjugación en
$\pi_1^{\mathrm{mod}}(\check{\mathbb P}^{r}-\check X,\bar\eta)$ de los
homomorfismos compuestos de inercia
\begin{center}
\begin{tikzcd}[column sep=huge,row sep=huge]
\mathbb Z_\ell(1) \arrow[r]
  & \pi_1^{\mathrm{mod}}(D-D\cap\check X,\bar\eta) \arrow[d] & {}\\
{}
  & \pi_1^{\mathrm{mod}}(\check{\mathbb P}^{r}-\check X,\bar\eta) & .
\end{tikzcd}
\end{center}
Por consiguiente, los ciclos evanescentes determinados por estos
homomorfismos son conjugados salvo signo por elementos de
$\pi_1^{\mathrm{mod}}(\check{\mathbb P}^{r}-\check X,\bar\eta)$.
La demostración termina por 6.1.1, que da la sobreyección
\[
\pi_1^{\mathrm{mod}}(D-D\cap\check X)
\longrightarrow
\pi_1^{\mathrm{mod}}(\check{\mathbb P}^{r}-\check X).
\]

\pageoriginale

\subsubsection*{Corolario 6.7.}

Bajo las hipótesis de 6.6, supongamos que $X$ satisface la condición
(LV). Entonces:

\medskip
\noindent
\textup{6.7.1.} La representación de
$\pi_1(D-D\cap\check X)$ sobre
$E^{n-1}(X_{\bar\eta},\mathbb Q_\ell)$ es absolutamente irreducible.

\medskip
\noindent
\textup{6.7.2.} Si $n$ es par, todo subgrupo de índice finito de
$\pi_1(D-D\cap\check X)$ actúa de manera absolutamente irreducible
sobre $E^{n-1}(X_{\bar\eta},\mathbb Q_\ell)$.

\subsubsection*{Demostración.}

Podemos reemplazar el módulo considerado por su giro de Tate
\[
E=E^{n-1}(X_{\bar\eta},\mathbb Q_\ell(m)).
\]
Sea $\bar E=E\otimes_{\mathbb Q_\ell}\bar{\mathbb Q}_\ell$, y elijamos
$0\ne\chi\in\bar E$. Para probar 6.7.1 debemos mostrar que el subespacio
$\bar{\mathbb Q}_\ell$-lineal de $\bar E$ engendrado por $\chi$ y sus
conjugados bajo $\pi$ es todo $\bar E$. Por 6.6.1 y 6.6.2 basta probar
que este subespacio contiene uno de los ciclos evanescentes $\delta_i$.

Como $X$ satisface (LV), la restricción de la forma del producto cup a
\[
H^{n-1}(X,\mathbb Q_\ell(m))
\lhook\joinrel\longrightarrow
H^{n-1}(X_{\bar\eta},\mathbb Q_\ell(m))
\]
es no degenerada; por tanto, también lo es su restricción al ortogonal
$E^{n-1}(X_{\bar\eta},\mathbb Q_\ell(m))$. Por 6.6.1 existe, pues, al
menos un $i$ tal que
\[
\chi\cdot\delta_i\ne0.
\]
Para este $i$, 6.5 proporciona un elemento $\sigma_i$ del grupo de
inercia en $s_i$ tal que
\begin{equation}\tag{6.7.3}
\sigma_i(\chi)-\chi
=\text{un múltiplo no nulo de }\delta_i.
\end{equation}
Esto prueba 6.7.1.

\pageoriginale

Para 6.7.2, $n-1$ es impar, luego
$\delta_i\cdot\delta_i=0$. Al iterar 6.5.2 se obtiene, para todo $k$,
\begin{equation}\tag{6.7.4}
\sigma_i^k(\chi)
=\sisi{X}{\chi}
+k\,\sigma_i(\sisi{x}{\chi}\cdot\delta_i)\delta_i.
\end{equation}
Así, un subespacio $\mathbb Q_\ell$-lineal $M$ de $E$ es estable bajo
la acción de $\pi$ si y solo si lo es bajo un subgrupo de índice finito:
ambas condiciones equivalen a
\[
\chi\in M
\quad\Longrightarrow\quad
\delta_i\in M\quad\text{para todo $i$ tal que }
\chi\cdot\delta_i\ne0.
\]
La afirmación se reduce, por tanto, a 6.7.1.

\subsection*{Bibliografía}

\begin{description}
\item[{[1]}] P. Deligne, \emph{Théorème de Lefschetz et critères de
dégénérescence de suites spectrales}, Publ. Math. IHÉS \textbf{35}
(1968), 259--278.

\item[{[2]}] S. Kleiman, \emph{Algebraic cycles and the Weil
conjectures}, en \emph{Dix exposés sur la cohomologie des schémas},
North-Holland, 1968.

\item[{[3]}] S. Lefschetz, \emph{L'analyse situs et la géométrie
algébrique}, Gauthier-Villars, París, 1924.

\item[{[4]}] D. Lieberman, \emph{Higher Picard Varieties}, Amer. J.
Math. \textbf{\sisi{15}{90}} (4) (1968), 1165--1199.

\item[{[5]}] D. Mumford, inédito.

\item[{[6]}] A. Weil, \emph{Variétés kählériennes}, Hermann, París,
1958.
\end{description}
